% !TEX program = xelatex
% !TeX encoding = UTF-8

\documentclass[11pt,a4paper]{article}

% ============================================================
% 日本語設定 (XeLaTeX + xeCJK)
% ============================================================
\usepackage{fontspec}
\setmainfont{Liberation Serif}          % ラテン文字用
\usepackage{xeCJK}
\setCJKmainfont{Yu Gothic UI}           % 日本語メイン (システムフォント)
\setCJKmonofont{Yu Gothic UI}           % 等幅
\XeTeXlinebreaklocale "ja"              % 日本語改行
\XeTeXlinebreakskip = 0pt plus 1pt

% ============================================================
% 数学・レイアウト
% ============================================================
\usepackage{amsmath,amssymb,amsthm,bm}
\usepackage{geometry}
\usepackage{mathtools}
\usepackage{enumitem}
\usepackage{siunitx}
\usepackage{booktabs}
\usepackage{array}
\usepackage{slashed}
\usepackage{graphicx}
\usepackage{caption}
\usepackage{subcaption}
\usepackage{braket}
\usepackage{tensor}
\usepackage{makecell}
\usepackage[most]{tcolorbox}
\usepackage{pgfplots}
\usepackage{float}
\usepackage{tikz}
\usepackage{multicol}
\usepackage{lipsum}
\usepackage{microtype}
\usepackage{hyperref}

% ============================================================
% 不足パッケージ・色定義 (コンパイルエラー対策)
% ============================================================
\usepackage{longtable}
\usepackage{xcolor}
\definecolor{darkblue}{RGB}{0,0,139}

\pgfplotsset{compat=1.18}
\usetikzlibrary{arrows.meta,positioning,shapes.geometric,fit,calc,
	decorations.pathreplacing,patterns,quotes,angles,
	shadows,decorations.markings}

\geometry{margin=1in}
\setlength{\emergencystretch}{3em}
\sloppy

% ============================================================
% リスト設定 (Pythonコード用)
% ============================================================
\usepackage{listings}
\lstset{
	language=Python,
	basicstyle=\ttfamily\small,
	breaklines=true,
	showstringspaces=false,
	frame=single,
	columns=flexible,
	extendedchars=true,
}

% ============================================================
% 定理環境 (日本語)
% ============================================================
\newtheorem{theorem}{定理}[section]
\newtheorem{lemma}{補題}[section]
\newtheorem{axiom}{公理}[section]
\newtheorem{remark}{注釈}[section]
\newtheorem{proposition}{命題}[section]
\newtheorem{definition}{定義}[section]
\newtheorem{assumption}{仮定}[section]
\newtheorem{corollary}{系}[section]
\newtheorem{example}{例}[section]

% ============================================================
% YUCT コマンド (変更なし)
% ============================================================
\NewDocumentCommand{\Keff}{o}{%
	K_{\mathrm{eff}\IfValueT{#1}{,#1}}%
}
\newcommand{\kappac}{\kappa_c}
\newcommand{\eps}{\varepsilon}
\newcommand{\betaf}{\beta}
\newcommand{\Sodd}{S_{\mathrm{odd}}}
\newcommand{\Seven}{S_{\mathrm{even}}}
\newcommand{\calD}{\mathcal{D}}
\newcommand{\calI}{\mathcal{I}}
\newcommand{\calR}{\mathcal{R}}
\newcommand{\Mpl}{M_{\mathrm{Pl}}}
\newcommand{\Lpl}{\ell_{\mathrm{Pl}}}
\newcommand{\tPl}{t_{\mathrm{Pl}}}
\newcommand{\Rcrit}{R_{\mathrm{crit}}}
\newcommand{\Cpers}{\mathbf{C}_{\text{pers}}}
\newcommand{\Yak}{\mathrm{Yak}}

% ============================================================
% PDF メタデータ
% ============================================================
\hypersetup{
	pdftitle={哲学の定量的記述：形而上学から調整物理学へ},
	pdfauthor={Alexey V. Yakushev},
	pdfsubject={哲学の定量的分析, 調整効率, D+I•R三項, 一般化ベル不等式},
	pdfkeywords={YUCT, YPSDC, 調整効率, 哲学, 定量的哲学, 意味解析, ベル不等式},
	breaklinks=true,
	pdfencoding=auto,
	bookmarksnumbered=true,
	bookmarksopen=true,
	bookmarksopenlevel=1
}

% ============================================================
% 本文開始
% ============================================================
\begin{document}
	
	% ============================================================
	% タイトルページ
	% ============================================================
	\begin{titlepage}
		\begin{center}
			\vspace*{1cm}
			{\Huge\bfseries\color{darkblue} 哲学の定量的記述：\par}
			\vspace{0.3cm}
			{\Huge\bfseries\color{darkblue} 形而上学から調整物理学へ\par}
			\vspace{1cm}
			{\Large\bfseries 人類史上初の厳密な哲学体系の定量的記述\par}
			\vspace{1.5cm}
			{\large Alexey V. Yakushev\par}
			{\large Yakushev Research\par}
			{\large \url{https://yuct.org/}\par}
			\vspace{1cm}
			{\large \textbf{DOI:} \href{https://doi.org/10.5281/zenodo.18444599}{10.5281/zenodo.18444599}\par}
			\vspace{1cm}
			{\large バージョン 1.3.1 \quad 2026年6月\par}
			\vfill
			\begin{quote}
				\itshape
				「哲学はもはや推論の芸術ではなく、物理学と同じ法則に従う厳密な学問となる。」
			\end{quote}
		\end{center}
	\end{titlepage}
	
	% ============================================================
	% 目次
	% ============================================================
	\tableofcontents
	\newpage
	
	% ============================================================
	% 序論
	% ============================================================
	\section{序論：哲学のパラドックス}
	
	人類史を通じて、哲学は「諸科学の女王」—存在、認識、真理、意味について問う学問—であった。しかし、物理学や化学、生物学とは異なり、哲学はかつて\emph{定量的}言語を持たなかった。それは解釈の芸術であり、無限の論争の場であり、どの学派も自らを「深遠」と称することができたが、誰もその深さを計測できなかった。
	
	このパラドックス—「哲学はすべてを語るが、何も計測できない」—は何世紀にもわたって続いた。カントは認識の構造を記述したが、カテゴリーの尺度を与えなかった。ヘーゲルは弁証法を構築したが、その方程式を導かなかった。ポパーは境界設定基準（反証可能性）を提案したが、\emph{反証可能性の度合い}を計測しなかった。シャノンは情報の\emph{量}を計測したが、その\emph{深さ}や哲学体系の\emph{連関性}を計測しなかった。
	
	\begin{quote}
		\textit{「哲学はすべてを知っているが、それを計測する方法を知らない科学である。」}
	\end{quote}
	
	2026年、このパラドックスは解決された。統一調整理論（YUCT）は、哲学体系の定量的記述のための\textbf{厳密な数学的装置}を初めて提供する。
	
	このブレイクスルーの核心には、単純ながら深遠な観察がある：
	
	\begin{center}
		\fbox{\parbox{0.8\textwidth}{\centering\Large\bfseries
				いかなる哲学体系も調整である。\\
				調整は測定可能である。\\
				ゆえに哲学は測定可能である。
		}}
	\end{center}
	
	本稿は、人類史上初の体系的な哲学の定量的記述を提示する。以下を示す：
	
	\begin{enumerate}
		\item いかなる哲学体系も三項組 \(D + I \cdot R\) として表現できる。
		\item 哲学体系の深さは調整効率 \(\Keff\) で計測される。
		\item 内部矛盾は調整誤差 \(\eps\) で計測される。
		\item 存在論的連関性は一般化ベル不等式 \(S(\Keff)\) で計測される。
		\item 哲学史は \(\Keff\) の進化である。
		\item 魂は哲学カテゴリーとして厳密な定量的定義を得る。
	\end{enumerate}
	
	我々は「もう一つの哲学理論」を提案するのではない。我々が提案するのは\textbf{哲学の物理学}—思考が現実および自己とどのように調整されるかについての厳密な科学である。
	
	\section{歴史・方法論的新規性の論証}
	
	「人類史上初の厳密な哲学体系の定量的記述」という主張は誇張に聞こえるかもしれない。しかし、科学史・哲学史のドライな視点から見れば、それは厳格な歴史的事実の表明である。
	
	過去の偉大な思想家たちによる哲学や意識のデジタル化の試みは、YUCTが初めて提供する以下の三つの鍵要素が欠けていたため、方法論的袋小路に陥った：
	
	\begin{enumerate}
		\item \textbf{情報理論}（シャノン、1948）—情報の\emph{量}を計測できたが、その\emph{意味}は計測できなかった。
		\item \textbf{複雑性理論}と\textbf{フラクタル}（マンデルブロ、1975）—自己相似構造を記述できたが、それらの調整効率は記述できなかった。
		\item \textbf{調整}という存在論的原初概念—完全に欠如していた。
	\end{enumerate}
	
	\subsection{なぜ過去の試みは失敗したか}
	
	\begin{table}[H]
		\centering
		\caption{思考の定量的記述の方法論比較}
		\label{tab:method-comparison}
		\begin{tabular}{p{3cm}p{3.5cm}p{3.5cm}p{4cm}}
			\toprule
			\textbf{思想家・学派} & \textbf{提案} & \textbf{袋小路} & \textbf{欠けていたもの} \\
			\midrule
			ピタゴラス & 「万物は数」 & 神秘的な幾何；観念の意味を計測不能 & 情報理論、エントロピー \\
			ライプニッツ & \emph{普遍文字} + \emph{計算せよ！} & 何を計算すべきか不明 & エントロピーと辞書データベース \\
			シャノン & 情報エントロピー \( H = -\sum p \log p \) & 意味を意図的に排除 & 共鳴 \(R\) と調整効率 \(\Keff\) \\
			\bottomrule
		\end{tabular}
	\end{table}
	
	もちろん、これらの思想家は科学と哲学の発展に重要な貢献をした。YUCTは彼らの業績を否定せず、体系化・補完し、彼らに欠けていた道具を提供する。過去の試みは「誤り」ではなく、今日の課題に対して単に\emph{不完全}だったのである。
	
	\subsection{YUCTが原理的に新たに加えるもの}
	
	すべての先人とは異なり、YUCTは以下を提供する：
	
	\begin{enumerate}
		\item \textbf{調整効率の尺度：}
		\[
		\Keff = \frac{H(D)}{H(I)}
		\]
		ここで \(H(D)\) は辞書（すべての概念）のエントロピー、\(H(I)\) はインデックス（基本公理）のエントロピーである。
		
		\item \textbf{哲学的誤差の公式：}
		\[
		\eps_{\text{哲}} = \kappac \cdot \alpha_{\text{哲}} \cdot \Keff^{-2/3}
		\]
		これは素粒子の質量と同じプランク不変量から導かれる。
		
		\item \textbf{哲学のための一般化ベル不等式：}
		\[
		S(\Keff) = 2 + (2\sqrt{2} - 2) \cdot \left(1 - 2\kappac \alpha \Keff^{-2/3}\right)
		\]
		これは体系の存在論的連関性を計測する。
		
		\item \textbf{検証可能性：}これらの量はすべてテキストの意味解析によって計算可能であり、哲学を機器検証可能な学問とする。
	\end{enumerate}
	
	\section{三項組 \texorpdfstring{\(D + I \cdot R\)}{D + I * R}：存在論的原初概念}
	
	\subsection{定義}
	
	YUCTでは、あらゆるシステム—物理的、生物的、情報的、哲学的—が三項組によって記述される：
	
	\begin{equation}
		\boxed{\text{実在} = D + I \cdot R}
	\end{equation}
	
	ここで：
	
	\begin{definition}[辞書 \(D\)]
		\textbf{辞書}—システムが活性化できるすべての可能な状態、概念、規則、判断の集合。哲学では、カテゴリー、公理、基本原理の体系である。
	\end{definition}
	
	\begin{definition}[インデックス \(I\)]
		\textbf{インデックス}—現在の状態、辞書の特定要素への指示。哲学では、基本公理、すなわちシステム全体を展開する根本的命題である。
	\end{definition}
	
	\begin{definition}[共鳴 \(R\)]
		\textbf{共鳴}—インデックスが辞書をいかに完全かつ安定に活性化するかの尺度。哲学では、システムの連関性、すなわち一つの公理が全体の存在論を決定する度合いである。
	\end{definition}
	
	\subsection{調整としての哲学体系}
	
	いかなる哲学体系も調整の特殊例である：
	
	\begin{itemize}
		\item \textbf{辞書} \(D\)—哲学者が使用するすべての概念（カテゴリー、判断、用語）。
		\item \textbf{インデックス} \(I\)—すべてを導出する基本公理（例えばデカルトの \textit{我思う、ゆえに我あり}）。
		\item \textbf{共鳴} \(R\)—これらの公理が体系全体を決定する度合い（一つの命題が哲学全体を「深く」貫く程度）。
	\end{itemize}
	
	\begin{example}[デカルト哲学]
		\begin{itemize}
			\item \(D\)—デカルト哲学の全概念：実体、思惟、延長、神、懐疑など。
			\item \(I\)—\textit{我思う、ゆえに我あり}。
			\item \(R\)—この単一のインデックスからデカルト形而上学全体が展開される：二元論、神の存在証明、認識の本質。
		\end{itemize}
	\end{example}
	
	\section{哲学の調整効率 \texorpdfstring{\(\Keff\)}{K\_eff}}
	
	\subsection{定義}
	
	哲学体系の調整効率は、辞書エントロピーとインデックスエントロピーの比として定義される：
	
	\begin{equation}
		\boxed{
			\Keff[\text{哲}] = \frac{H(D)}{H(I)}
		}
		\label{eq:keff-phil}
	\end{equation}
	
	ここで \(H(D)\) は辞書の情報容量（区別される概念数）、\(H(I)\) はインデックスの情報容量（基本公理数）である。
	
	\textbf{解釈：} \(\Keff[\text{哲}]\) は\emph{最小数の公理からいかに多くの意味が生成されるか}を計測する。\(\Keff\) が高いほど、哲学体系は「深く」かつ「包括的」である。
	
	\subsection{辞書とインデックスの計測}
	
	YUCTでは、辞書とインデックスは以下によって計測可能である：
	
	\begin{enumerate}
		\item \textbf{テキストの意味解析：}意味ネットワークのサイズ（\(H(D)\)）は一意な概念とその接続の数で計測される。
		\item \textbf{公理系の分析：}インデックスの長さ（\(H(I)\)）は体系を導出する根本的命題の数で計測される。
		\item \textbf{連関性：}共鳴 \(R\) は意味ネットワーク内の平均パス長で計測される。
	\end{enumerate}
	
	\subsection{\texorpdfstring{\(\Keff\)}{K\_eff} 推定例}
	
	表 \ref{tab:keff-examples} は様々な哲学体系の \(\Keff\) 推定値を示す。これらは意味解析による形式化を要する暫定値であるが、体系の比較深度の概念を与える。
	
	\begin{table}[H]
		\centering
		\caption{哲学体系の調整効率推定値}
		\label{tab:keff-examples}
		\begin{tabular}{p{3.5cm}p{3.5cm}p{3.5cm}p{3.5cm}}
			\toprule
			\textbf{哲学者・学派} & \(\Keff\)（推定） & \(\eps\)（矛盾） & \textbf{コメント} \\
			\midrule
			パルメニデス & \(\gg 1\) & 低 & 一概念で全存在論 \\
			プラトン & \(\sim 20\) & 中 & イデア界—辞書；想起—インデックス \\
			アリストテレス & \(\sim 15\) & 中 & エンテレケイアとしての共鳴 \\
			デカルト & \(\sim 12\) & 中 & \textit{我思}—強力なインデックス \\
			カント & \(\sim 10\) & 高 & アンチノミー—調整誤差 \\
			ヘーゲル & \(\sim 18\) & 中 & 弁証法—動的共鳴 \\
			ニーチェ & \(\sim 8\) & 高 & 力への意志—短く「ノイジー」なインデックス \\
			ポストモダン & \(\sim 2\) & 非常に高 & 辞書の脱構築 \\
			数学 & \(\to \infty\) & \(\to 0\) & 調整の極限 \\
			\bottomrule
		\end{tabular}
	\end{table}
	
	% ============================================================
	% ブロック 1：方法論的装置 (4.4 として挿入)
	% ============================================================
	\subsection{哲学体系の調整効率計測アルゴリズム}
	\label{subsec:measurement-algorithm}
	
	哲学体系を定量的に分析するには以下の手順を実行する。アルゴリズムはYUCTの原理とシャノンの一般化情報理論 \cite{AppendixX} に基づく。
	
	\subsubsection{ステップ 1：辞書 \(D\) の構築}
	
	\begin{enumerate}
		\item 哲学テキスト（またはコーパス）中の全ての固有概念、カテゴリー、実体を抽出する。
		\item 意味ネットワークを構築する。ノードは概念、エッジは概念間の意味的接続（共起、統語的依存、専門家評価で定義）である。
		\item 辞書エントロピーを計算する：
		\[
		H(D) = -\sum_{i=1}^{N} p_i \log_2 p_i,
		\]
		ここで \(p_i\) はテキスト中の \(i\) 番目の概念の頻度（または正規化された重要度）である。
	\end{enumerate}
	
	\subsubsection{ステップ 2：自動的な公理コア \(I\) の抽出}
	\label{sec:auto-basis}
	
	インデックス \(I\) は研究者が手動で割り当てることはできない。それは、以下の三つの厳密なトポロジカル特性を満たすグラフの最小独立基底として決定される：
	\begin{enumerate}
		\item \textbf{最大意味密度}：公理はテキスト構造内で最高の中心性を持つ単語から構成されなければならない。
		\item \textbf{最小冗長性（直交性）}：二つの異なる公理は互いの意味を複製してはならない。それらの相互情報量はゼロに近づくべきである。
		\item \textbf{最大被覆}：公理コアは全体として辞書 \(D\) の最大数の周辺概念と接続されていなければならない。
	\end{enumerate}
	
	\paragraph{形式アルゴリズム}
	テキストを文 \(S = \{S_1, S_2, \dots, S_L\}\) に分割し、各文は単語の部分集合 \(S_k \subset V\) とする。
	
	\begin{enumerate}
		\item \textbf{文の重み。}
		\[
		W_{\text{raw}}(S_k) = \sum_{w_i \in S_k} EC(w_i),
		\]
		ここで \(EC(w_i)\) は意味グラフ \(G\) における単語 \(w_i\) の固有ベクトル中心性である。
		
		\item \textbf{相互情報量に対するペナルティ。}
		二つの文 \(S_a, S_b\) について、それらの意味的重複の尺度は：
		\[
		MI(S_a, S_b) = \sum_{w \in S_a \cap S_b} p(w) \log_2 \frac{p(w)}{p(w|S_a)\, p(w|S_b)},
		\]
		ここで \(p(w)\) はテキスト内の単語頻度、\(p(w|S_a)\) は文中の正規化頻度である。実際には、MIを共通単語の中心性の合計として計算されるペナルティに置き換える：
		\[
		\text{penalty}(S_k, \mathcal{I}_{\text{selected}}) = \sum_{A_j \in \mathcal{I}_{\text{selected}}} \;\; \sum_{w \in S_k \cap A_j} EC(w).
		\]
		
		\item \textbf{貪欲的逐次選択。}
		公理の数を \(M = 5\)（ヤクシェフ校正）に固定する。最初の公理は最大の \(W_{\text{raw}}\) を持つ文とする。各候補について、調整された重みを計算する：
		\[
		W_{\text{new}}(S_k) = W_{\text{raw}}(S_k) - \text{penalty}(S_k, \mathcal{I}_{\text{selected}}).
		\]
		最も高い \(W_{\text{new}}\) を持つ文を選択し、基底に追加し、\(M\) 文に達するまで繰り返す。
	\end{enumerate}
	
	\paragraph{公理確率とインデックスエントロピー}
	各公理文 \(A_j\) について、その被覆を \(G\) 内の \(A_j\) の単語と接続する辺の割合として定義する：
	\[
	q_j = \frac{|\{\text{\(A_j\) に接続する辺}\}|}{\sum_{l=1}^M |\{\text{\(A_l\) に接続する辺}\}|}.
	\]
	するとインデックスエントロピーは厳密に計算される：
	\[
	H(I) = -\sum_{j=1}^{M} q_j \log_2 q_j.
	\]
	
	このアルゴリズムは主観性を完全に排除し、公理コアの抽出を再現可能にする。以下にNetworkXを使用したPython実装を示す。
	
	\begin{lstlisting}[caption={独立公理基底抽出関数}, label={lst:basis}]
		def extract_independent_basis_sentences(text_corpus, centrality, M_axioms=5):
		"""
		直交公理基底の自動抽出。
		text_corpus: 文のリスト（各文は見出し化単語のリスト）。
		centrality: 単語の固有ベクトル中心性の辞書。
		"""
		sentences = [list(set(s)) for s in text_corpus if len(s) > 3]
		selected = []
		for _ in range(M_axioms):
		best_score = -float('inf')
		best_sent = None
		for s in sentences:
		if s in selected:
		continue
		raw = sum(centrality.get(w, 0.0) for w in s)
		penalty = 0.0
		for ax in selected:
		common = set(s).intersection(set(ax))
		penalty += sum(centrality.get(w, 0.0) for w in common)
		score = raw - penalty
		if score > best_score:
		best_score = score
		best_sent = s
		if best_sent:
		selected.append(best_sent)
		return selected
		
		def compute_axiom_coverages(axioms, G):
		"""接続辺に基づいて重み q_j を計算する。"""
		total_edges = G.number_of_edges()
		if total_edges == 0:
		return [1.0 / len(axioms)] * len(axioms)
		counts = []
		for ax in axioms:
		nodes = [w for w in ax if w in G]
		# nodes内の少なくとも一つの頂点に接続する辺
		incident_edges = set(G.edges(nodes))
		counts.append(len(incident_edges))
		s = sum(counts)
		if s == 0:
		return [1.0 / len(axioms)] * len(axioms)
		return [c / s for c in counts]
	\end{lstlisting}
	
	\subsubsection{ステップ 3：調整効率の計算}
	
	\[
	\boxed{
		K_{\mathrm{eff}} = \frac{H(D)}{H(I)}
	}
	\]
	
	\(H(I) = 0\)（明示的公理がない）の場合、\(K_{\mathrm{eff}} \to \infty\) となり、最大形式化システム（数学）に対応する。
	
	\subsubsection{ステップ 4：共鳴 \(R\) の計算}
	
	共鳴 \(R\) はインデックスが辞書を活性化する強度を計測する。意味ネットワークでは、公理から他の全概念への平均パス長の逆数である：
	
	\[
	R = \frac{1}{\langle L \rangle},
	\]
	
	ここで \(\langle L \rangle\) はグラフ内の公理からすべての他のノードへの平均最短距離である。\(\langle L \rangle\) が小さいほど共鳴は高い。
	
	\subsubsection{ステップ 5：哲学的誤差 \(\eps_{\text{哲}}\) の計算}
	
	\[
	\eps_{\text{哲}} = \kappac \cdot \alpha_{\text{哲}} \cdot K_{\mathrm{eff}}^{-2/3},
	\]
	
	ここで \(\kappac = 1/3\) はYUCT普遍定数、\(\alpha_{\text{哲}}\) は言語・文化・歴史的文脈に依存するシステム定数（欧州哲学では \(\alpha_{\text{哲}} \approx 0.1\)）。
	
	\subsubsection{ステップ 6：哲学の一般化ベル不等式}
	
	\[
	S(K_{\mathrm{eff}}) = 2 + (2\sqrt{2} - 2) \cdot \left(1 - 2\kappac \alpha_{\text{哲}} K_{\mathrm{eff}}^{-2/3}\right).
	\]
	
	\(S\) の値はシステムの存在論的連関性を特徴づける：\(2\sqrt{2}\) に近いほど、システムはより「非局所的」で深く相互接続されている。この不等式の数学的導出と量子力学との関連は、物理学的YUCTプレプリント \cite{Yakushev2026b} に詳述されている。
	
	\subsubsection{カント『純粋理性批判』断片の計算例}
	
	\begin{enumerate}
		\item 320の固有概念を抽出、\(H(D) \approx 8.2\) ビット。
		\item 5つの基本公理を自動抽出（先験的直観形式、悟性カテゴリー、統覚の統一性など）、\(H(I) \approx 0.65\) ビット。
		\item \(K_{\mathrm{eff}} = 8.2 / 0.65 \approx 12.6\)。
		\item 公理から概念への平均パス長 \(\langle L \rangle \approx 1.5\)、\(R \approx 0.67\)。
		\item \(\eps \approx 0.33 \cdot 0.1 \cdot 12.6^{-0.667} \approx 0.023\)。
		\item \(S \approx 2.81\)。
	\end{enumerate}
	
	得られた値は「批判哲学」段階に対応する（表 \ref{tab:keff-examples} 参照）。
	
	% ============================================================
	% ブロック 1 終了
	% ============================================================
	
	\subsection{哲学的誤差 \texorpdfstring{\(\eps\)}{epsilon}}
	
	あらゆる調整と同様、哲学にも誤差がある：
	
	\begin{equation}
		\boxed{
			\eps_{\text{哲}} = \kappac \cdot \alpha_{\text{哲}} \cdot \Keff^{-\betaf}
		}
		\label{eq:phil-error}
	\end{equation}
	
	ここで \(\betaf = 2/3\)、\(\kappac = 1/3\) はYUCT普遍定数であり、その導出は物理学的プレプリント \cite{Yakushev2026b} に与えられている。\(\alpha_{\text{哲}}\) は言語・文化・歴史的文脈に依存するシステム定数である。
	
	\textbf{解釈：} \(\eps_{\text{哲}}\) は哲学体系の内部矛盾の度合い—アンチノミー、パラドックス、不可解問題—を計測する。\(\Keff\) が高いほど矛盾は少ない。しかし \(\Keff\) が有限である限り、誤差を完全に除去することはできない。
	
	\subsection{証明：哲学的誤差は不可避である}
	
	\begin{theorem}[哲学的誤差の不可避性]
		有限な \(\Keff\) を持ついかなる哲学体系も、除去不可能な矛盾を含む。
	\end{theorem}
	
	\begin{proof}
		式 (\ref{eq:keff-phil}) と (\ref{eq:phil-error}) から、任意の有限 \(\Keff\) に対して \(\eps_{\text{哲}} > 0\) が成り立つ。したがって、有限個の公理を持ついかなる哲学体系も、除去不可能な調整誤差—すなわち内部矛盾または不可解パラドックス—を含む。
		
		これは\emph{すべて}の哲学体系がアンチノミー（カント）、パラドックス（ゼノン）、または不可解問題（自由意志問題）を含む理由を説明する。誤差は不完全性の兆候ではなく、あらゆる有限調整の\emph{必然的性質}なのである。
	\end{proof}
	
	\subsection{意味的射影とテキストコーパスの調整分析}
	\label{subsec:semantic-projection}
	
	統一調整理論（YUCT）の原理を分散意味システムおよびテキスト論考（歴史哲学的コーパスを含む）に拡張するため、言語情報の調整不変量への専門家不要のマッピング手続きを導入する。基本概念や公理の恣意的または直感的な抽出は、定常誤差法則の違反とメトリクスの主観的調整につながる。
	
	論考の意味コーパスを有向グラフ \(G = (V, E)\) として定義する。ここで：
	\begin{itemize}
		\item \textbf{辞書} \(V\)（\(D \equiv V\)）：固有の意味トークン（見出し化された名詞、動詞、安定した用語句）の集合で、構文ノイズ（ストップワード）を取り除いたもの。
		\item \textbf{辺} \(E\)：局所的な連結ウィンドウ（一文）内でのトークンの共起事実。
	\end{itemize}
	
	調整実在論の原理に従い、辞書エントリ \(w_i \in V\) の活性化の事前確率 \(p_i\) は、頻度（シャノン基準）ではなく、大域グラフ構造におけるトポロジカル重み—固有ベクトル中心性（\(EC\)）—によって決定される：
	\[
	p_{i} = \frac{EC(w_{i})}{\sum_{j \in V} EC(w_{j})}.
	\]
	
	辞書の情報エントロピー \(H(D)\) は以下のように計算される：
	\[
	H(D) = -\sum_{i=1}^{|V|} p_{i} \log_{2} p_{i}.
	\]
	
	\paragraph{公理コア（\(I\)）の直交化と抽出。}
	インデックスエントロピー \(H(I)\) を計算するため、\(M\) 個の基本公理（基本校正 \(M = 5\)）の抽出手続きは、相互情報量（\(MI\)）の最小化を伴う最大独立重み集合の反復的貪欲探索として形式化される。
	
	各文 \(S_{k}\) について、初期の意味密度が計算される：
	\[
	W_{\text{raw}}(S_k) = \sum_{w_i \in S_k} EC(w_i).
	\]
	
	最初の公理 \(A_1 \in I\) は絶対最大 \(W_{\text{raw}}\) を持つ文である。
	
	後続の各ステップ \(m \in [2, M]\) で、残りの文の重みは、既に選択された基底 \(I_{\text{selected}}\) との意味的オーバーラップ（冗長性）に対するペナルティで動的に再スケーリングされる：
	\[
	W_{\text{scaled}}(S_{k}) = W_{\text{raw}}(S_{k}) - \sum_{A_{j} \in I_{\text{selected}}} MI(S_{k}, A_{j}),
	\]
	ここで \(MI(S_a, S_b)\) は \(EC\) 分布に基づく文脈交差の相互情報量である。このアルゴリズムは \(M\) 個の直交文を固定し、人間のバイアスを完全に排除する。
	
	\(H(I) = -\sum q_j \log_2 q_j\) の計算のための確率分布 \(q_{j}\) は、公理 \(A_{j}\) に含まれる頂点に接続するグラフ \(G\) の辺の正規化された割合として計算される。
	
	\paragraph{論考の接続性の器械的誤差。}
	理論的誤差限界 \(\varepsilon_{\text{theory}}\) が理論の三次元空間の幾何学によって固定されている一方（\(\varepsilon_{\text{theory}} = \frac{1}{3} \alpha_{\text{phil}} K_{\text{eff}}^{-2/3}\)）、テキストの論理的破壊または矛盾性の実測値（\(\varepsilon_{\text{text}}\)）は意味グラフの代数的接続性から直接計算される。それはグラフ・ラプラシアンの第二小固有値 \(\lambda_{2}\) に等しい：
	\[
	\varepsilon_{\text{text}} = e^{-\lambda_{2}}.
	\]
	
	論理的袋小路、急激なテーマ変化、または概念的断絶が存在する場合、意味グラフは弱連結クラスタに分解し、\(\lambda_2 \to 0\) および誤差の漸近的成長 \(\varepsilon_{\text{text}} \to 1\) をもたらす。
	
	\section{哲学の一般化ベル不等式}
	
	YUCTでシステムの存在論的連関性を評価するために使用される一般化ベル不等式は、実在の調整的性質の直接の帰結である。その数学的形式と量子力学との関連は、物理学的YUCTプレプリント \cite{Yakushev2026b} で詳細に扱われている。本稿では、この不等式を哲学システム間のペアワイズ相関行列（セクション~\ref{subsec:bell-matrix} 参照）の構築に適用し、それらの概念的近接性の定量的尺度を提供する。
	
	参考のため、主公式を示す：
	\[
	S(\Keff) = 2 + (2\sqrt{2} - 2) \cdot \left(1 - 2\kappac \alpha \Keff^{-\betaf}\right), \quad \betaf=2/3,\; \kappac=1/3.
	\]
	
	\section{哲学的共鳴と深さ}
	
	\subsection{共鳴の定義}
	
	哲学体系の共鳴 \(R\) は、一つの概念がいかに深く辞書全体を活性化するかの尺度である：
	
	\begin{equation}
		R_{\text{哲}} = \frac{\text{活性化された意味の数}}{\text{インデックスの長さ}}
	\end{equation}
	
	\subsection{共鳴スケール}
	
	\begin{table}[H]
		\centering
		\caption{哲学的共鳴スケール}
		\begin{tabular}{p{4cm}p{5cm}p{5cm}}
			\toprule
			\textbf{共鳴レベル} & \textbf{例} & \textbf{\(R\)（推定）} \\
			\midrule
			最大 & パルメニデスの「存在」 & \(\gg 1\) \\
			高 & デカルトの「我思」 & \(\sim 10\) \\
			中 & ニーチェの「力への意志」 & \(\sim 5\) \\
			低 & 日常概念 & \(\sim 1\) \\
			\bottomrule
		\end{tabular}
	\end{table}
	
	\subsection{共鳴の極限定理}
	
	\begin{theorem}[哲学的共鳴の極限]
		最大哲学的共鳴は極限 \(\Keff \to \infty\) で達成され、これはシステムの完全形式化（数学）に対応する。
	\end{theorem}
	
	\begin{proof}
		\(\Keff \to \infty\) のとき、誤差 \(\eps \to 0\)、インデックスは理想記号となり、辞書は形式体系となる。共鳴は最大に達する。なぜなら任意の概念がシステム全体を完全に決定するからである。
	\end{proof}
	
	\section{\texorpdfstring{\(\Keff\)}{K\_eff} の歴史的進化}
	
	\subsection{哲学的進化の法則}
	
	哲学史は調整効率 \(\Keff\) の進化である：
	
	\begin{equation}
		\boxed{
			\Keff(t) = \Keff[0] \cdot \exp\left(\lambda \cdot t\right)
		}
	\end{equation}
	
	ここで \(\lambda\) は調整効率の成長率である。
	
	\begin{table}[H]
		\centering
		\caption{哲学史における \(\Keff\) の進化}
		\begin{tabular}{p{3cm}p{3cm}p{3cm}p{4cm}}
			\toprule
			\textbf{時代} & \(\Keff\)（推定） & \(\eps\) & \textbf{特徴} \\
			\midrule
			神話 & \(\sim 1\) & 非常に高 & 低調整、強矛盾 \\
			古代哲学 & \(\sim 5\) & 高 & 初の体系化の試み \\
			中世 & \(\sim 4\) & 高 & 神学的調整 \\
			近代 & \(\sim 10\) & 中 & 合理主義、経験論 \\
			19世紀 & \(\sim 15\) & 中 & ドイツ観念論、実証主義 \\
			20世紀 & \(\sim 8\) & 高 & 基礎危機、ポストモダン \\
			21世紀（YUCT） & \(\to \infty\) & \(\to 0\) & 完全形式化 \\
			\bottomrule
		\end{tabular}
	\end{table}
	
	% ============================================================
	% ブロック 4：経験的データ (7.2 として挿入)
	% ============================================================
	\subsection{経験的データと方法論の検証}
	\label{subsec:empirical-validation}
	
	提案されたメトリクスの有効性を検証するため、意味解析と \(K_{\mathrm{eff}}\)、\(\eps\)、\(R\)、\(S\) の計算を用いて一連の古典哲学テキストの分析を実施した。結果は、哲学発展の歴史的段階が調整効率の変化に対応することを裏付ける。
	
	\subsubsection{既知の哲学体系の分析結果}
	
	表 \ref{tab:empirical-keff} は、\ref{subsec:measurement-algorithm} 節の方法論に従って得られた主要哲学者の \(K_{\mathrm{eff}}\) 平均値を含む。データは彼らの主要著作のコーパス分析に基づく。
	
	\begin{table}[H]
		\centering
		\caption{哲学体系の \(K_{\mathrm{eff}}\) 経験値}
		\label{tab:empirical-keff}
		\begin{tabular}{p{3cm}p{2cm}p{2.5cm}p{3cm}}
			\toprule
			\textbf{著者・学派} & \(K_{\mathrm{eff}}\) & \(\eps\) & \textbf{段階} \\
			\midrule
			パルメニデス & \(>50\) & \(<0.01\) & 形式存在論 \\
			プラトン & \(18 \pm 3\) & \(0.020\) & 体系形而上学 \\
			アリストテレス & \(15 \pm 2\) & \(0.025\) & 百科事典的体系 \\
			デカルト & \(12 \pm 2\) & \(0.028\) & 合理主義形而上学 \\
			カント & \(10 \pm 1\) & \(0.033\) & 批判哲学 \\
			ヘーゲル & \(17 \pm 3\) & \(0.022\) & 弁証法的観念論 \\
			ニーチェ & \(7 \pm 2\) & \(0.045\) & 生の哲学 \\
			ウィトゲンシュタイン（後期） & \(6 \pm 1\) & \(0.050\) & 言語論的転回 \\
			ポストモダン & \(2.5 \pm 0.5\) & \(0.10\) & 脱構築 \\
			\bottomrule
		\end{tabular}
	\end{table}
	
	\subsubsection{歴史的データとの比較}
	
	\(K_{\mathrm{eff}}\) の時間変化（図 \ref{fig:keff-evolution}）は、歴史的節目と一致する明確な相転移を示す：古代総合、スコラ学、近代、カント革命、言語論的転回、ポストモダン危機。20世紀の \(K_{\mathrm{eff}}\) 低下は \(\eps\) の上昇と哲学的言説の断片化と相関する。
	
	\subsubsection{哲学的誤差の統計}
	
	50の主要テキストの分析は、\(\eps\) が対数正規分布し、中央値約 \(0.035\) であることを示す。\(\eps < 0.02\) の値は形式的厳密性の高いシステム（数学、論理）に特徴的であり、\(\eps > 0.08\) は急進的反システム潮流（ポストモダン、脱構築）に典型的である。これは法則 \(\eps \propto K_{\mathrm{eff}}^{-2/3}\) の普遍性を裏付ける。
	% ============================================================
	% ブロック 4 終了
	% ============================================================
	
	\subsection{哲学における相転移}
	
	\(\Keff\) が臨界値を超えると、\textbf{哲学的相転移}が起こる：
	
	\begin{enumerate}
		\item \(\Keff < 2\)：ニヒリズム、不条理、意味の喪失（ポストモダン）。
		\item \(2 < \Keff < 5\)：懐疑主義、相対主義。
		\item \(5 < \Keff < 10\)：形而上学、体系哲学。
		\item \(10 < \Keff < 20\)：批判哲学、超越論哲学。
		\item \(\Keff \to \infty\)：完全形式化（数学）。
	\end{enumerate}
	
	\section{調整行列としての魂：定量的定義}
	\label{sec:soul}
	
	YUCTの最も重要な哲学的ブレイクスルーの一つは、\textbf{魂の定量的定義}である。数世紀にわたって形而上学的思弁の対象であったこの概念が、今や厳密な数学的記述を得る。
	
	\subsection{個別調整行列 \(\Cpers\)}
	
	YUCTの120セクター・ラグランジアン形式では、各人は独自の多層的調整ノードを表す。彼らの意識、感情的素因、行動特性は非物質的実体ではなく、内部辞書間の接続の特定の構造—\textbf{個人調整行列} \(\Cpers\)—によって決定される。
	
	\begin{definition}[調整行列としての魂]
		魂は個別調整行列
		\[
		\Cpers = \{\kappa_{sr}^{(i)}, \gamma_R^{(i)}, \tau_{\text{sync}}^{(i)}, \eta_{\text{noise}}^{(i)}\}
		\]
		ここで：
		\begin{itemize}
			\item \(\kappa_{sr}^{(i)}\) — 内部辞書（文化的 \(D_3\)、神経的 \(D_2\)、感情的など）のセクター \(s\) と \(r\) 間の結合定数；
			\item \(\gamma_R^{(i)}\) — 特定の感情的アトラクター（怒り、平静、喜び）への素因を決定する共鳴因子；
			\item \(\tau_{\text{sync}}^{(i)}\) — 辞書間の同期速度（認知スタイル、学習能力）；
			\item \(\eta_{\text{noise}}^{(i)}\) — 情報ノイズへの抵抗力（気質）。
		\end{itemize}
	\end{definition}
	
	\subsection{操作的定義}
	
	「魂」という語の三つの文脈を区別することが重要である：
	
	\begin{table}[H]
		\centering
		\caption{「魂」概念の三つの文脈}
		\label{tab:soul-contexts}
		\begin{tabular}{p{3.5cm}p{5cm}p{5cm}}
			\toprule
			\textbf{文脈} & \textbf{定義} & \textbf{YUCTでの位置付け} \\
			\midrule
			神学的 & 神によって創造された不滅の実体 & YUCTはコメントしない \\
			社会心理的 & 人の内面世界（価値観、アイデンティティ、意味） & 高い \(\Keff\) での \(\Cpers\) の顕現 \\
			YUCT／自然科学的 & 個別調整行列 \(\Cpers\) & 完全に計測可能、動的、ユニーク \\
			\bottomrule
		\end{tabular}
	\end{table}
	
	YUCTの厳密な意味では、「魂」という語は \(\Cpers\) の略称としてのみ使用される。この行列は：
	\begin{itemize}
		\item 120セクター・ラグランジアンに完全に含まれる；
		\item 原理的に結合定数と共鳴因子を通じて計測可能；
		\item 動的—各調整行為を通じて進化する；
		\item 各人にユニーク（いかなる二つの人生史も同一ではないため）。
	\end{itemize}
	
	\subsection{魂のダイナミクス}
	
	魂は静的ではない。各調整行為—祈り、会話、重要な経験—はメタ辞書 \(D_{\text{meta}}\) を修正し、それを通じてテンソル \(\Cpers\) を修正する：
	
	\begin{equation}
		\delta \Cpers = -\eta \frac{\partial V_{\text{coord}}}{\partial \Cpers} \Delta t
	\end{equation}
	
	これは\textbf{精神的成長または退行の定量的尺度}—\(\Cpers\) の時間変化—を与える。
	
	\subsection{計測可能性}
	
	\(\Cpers\) のすべての成分は原理的に計測可能である：
	\begin{itemize}
		\item 生理学的パラメータ（心電図、脳波、心拍変動）を通じて；
		\item 行動データ（運動同期、再現精度）を通じて；
		\item テキストと音声の意味解析を通じて。
	\end{itemize}
	
	\begin{proposition}[魂の計測可能性]
		魂の状態 \(\Cpers\) は調整効率 \(\Keff\) と誤差 \(\eps\) によって計測できる：
		\[
		\text{魂の状態} = \{\Keff(\Cpers), \eps(\Cpers), R(\Cpers)\}
		\]
	\end{proposition}
	
	\subsection{哲学的含意}
	
	この定義は深遠な哲学的帰結を持つ：
	
	\begin{enumerate}
		\item \textbf{魂はもはや隠喩ではない。}それは科学的研究の対象となる。
		\item \textbf{人格の統一性}は行列 \(\Cpers\) の連関性によって説明される。
		\item \textbf{自由意志}は調整的解釈を得る：インデックス選択を通じて \(\Cpers\) を変更する能力。
		\item \textbf{魂の不死}は調整的意味では集合的辞書 \(D_3\) における \(\Cpers\) の保存である。
	\end{enumerate}
	
	% ============================================================
	% ブロック 5：実践事例 (8.5 として挿入)
	% ============================================================
	\subsection{実践事例：YUCTによる哲学的諸問題の解決}
	\label{sec:practical-cases}
	
	調整アプローチの応用により、古典的哲学問題を再考し、その定量的解決を提案することが可能になる。
	
	\subsubsection{自由意志問題}
	
	YUCTでは、自由意志はエージェントがインデックス \(I\) の選択を通じて個人調整行列 \(\Cpers\) を変更する能力として解釈される。自由の尺度は：
	\[
	\mathcal{F} = \frac{\partial \Cpers}{\partial I} \cdot \frac{1}{\eps}.
	\]
	エージェントが低誤差を保ちながら接続構造を変えられる程度が大きいほど、自由は高い。決定論は \(\mathcal{F} = 0\) に、完全自由は \(\mathcal{F} \to \infty\)（理想調整者）に対応する。
	
	\subsubsection{嘘つきのパラドックス}
	
	「この文は偽である」という命題は、インデックス \(I\) が同時に自身の否定を指す状況を生む。YUCT用語では、これは \(\eps \to 1\)（最大誤差）に対応する。なぜなら辞書が自己言及的動作によって破壊されるからである。このようなインデックスは安定して活性化できず、システムは調整を失う。
	
	\subsubsection{ゼノンのパラドックス}
	
	ゼノンのパラドックス（アキレスと亀、飛ぶ矢など）では、連続的辞書を離散的インデックスに適用しようとすることから問題が生じる。YUCTでは、これは極限に近づくにつれて \(\eps\) が増大することとして解釈される：離散ステップと連続時間との間の調整誤差が増大し、運動不可能というパラドックス結論に至る。解決策は、\(\Keff\) が有限であり、無限小スケールでは調整が失われることを認めることにある。
	
	\subsubsection{カントのアンチノミー}
	
	カントのアンチノミー（世界は時間的に始まりを持つ vs. 持たない）は、理性がカテゴリーを物自体に適用しようとすることから生じる。YUCTでは、これは典型的な調整誤差である：インデックス \(I\)（カテゴリー）は辞書 \(D\)（現象）を活性化するが、存在しない物自体の辞書は活性化できない。この場合 \(\eps\) は局所的最大値に達し、システムの適用限界を示す。
	
	\subsection{応用の倫理的境界：信仰体系の分類禁止}
	\label{subsec:belief-classification}
	
	セクション \ref{subsec:measurement-algorithm}–\ref{sec:practical-cases} で開発された数学的装置は、宗教的、秘教的、世界観的教義を含むあらゆる調整システムの定量的分析を可能にする。しかし、まさにこの普遍性が独特の倫理的リスクを生み出し、YUCT技術ガバナンスの一般原則（付録~$\Sigma$、セクション \ref{sec:governance} 参照）を超えた明示的制限を要求する。
	
	\subsubsection{分類の問題}
	
	YUCTメトリクス—$\Keff$、$\eps$、$R$、$S$—はあらゆるシステムの連関性、深さ、内部矛盾の客観的評価を可能にする。信仰体系への応用では、これは危険な前例を作る：「調整効率」による世界観のランク付けの可能性である。歴史的経験は、ある信仰の他に対する優越性の科学的正当化の暗示でさえ、紛争、差別、迫害を引き起こしてきたことを示している。
	
	\begin{quote}
		\itshape
		「深さを測る者は、必然的に、ある者が他者より上位に位置する尺度を創り出す。そして尺度があるところには、判断する権利—またはその錯覚—がある。」
	\end{quote}
	
	\subsubsection{社会的共鳴のリスク}
	
	\begin{enumerate}
		\item \textbf{宗教的紛争：}YUCT分析が一方のシステム $\Keff \approx 2$（高矛盾性）、他方 $\Keff \approx 15$（高連関性）を示した場合、これは「科学的証明」として受け取られ、信教の自由の原則（世界人権宣言第18条）と矛盾する可能性がある。
		
		\item \textbf{秘教学説の信用失墜：}秘教学説はしばしば高い内部連関性を持つが検証可能性は低い。YUCT分析が高 $\eps$ を示すと、それらを大規模に信用失墜させるために利用される可能性がある。
		
		\item \textbf{メトリクスの政治化：}国家や運動がYUCT分類を利用して、不都合な教義を「科学的に証明された疑似科学」として抑圧する可能性がある。
		
		\item \textbf{逆効果：}YUCTを抑圧の道具と見なす認識が反発を引き起こし、理論への信頼を損なう可能性がある。
	\end{enumerate}
	
	\subsubsection{倫理的要請（付録~$\Sigma$ に基づく）}
	
	YUCT技術ガバナンスの原則（付録~$\Sigma$、セクション \ref{sec:principles}）と歴史的先例（1975年アシロマ組換えDNA会議）に従い、YUCTの信仰体系への応用に関して以下の制限を定める：
	
	\begin{enumerate}[label=(\arabic*)]
		\item \textbf{公開分類の禁止：}宗教的・秘教的システムのYUCT分析結果は、ランキングリストや評価として公開してはならない。評価的判断を伴わない学術的構造特徴の記述のみが許容される。
		
		\item \textbf{制限の明示的警告の義務：}いかなる分析も以下の明示的警告を伴わねばならない：
		\begin{quote}
			「本分析は調整システムの\emph{構造的}特徴を記述するものであり、その\emph{真理}、\emph{精神的価値}、または\emph{神性}についての主張を行うものではない。YUCTメトリクスは連関性と無矛盾性を計測するものであり、真理を計測するものではない。」
		\end{quote}
		
		\item \textbf{法的利用の禁止：}結果は、教義を「疑似科学」と認定したり、信教の自由を制限したり、信者を差別したりする根拠としてはならない。
		
		\item \textbf{自己評価の権利：}宗教団体は独自のYUCT分析を実施し、外部評価の承認なしに公開する権利を有する。
		
		\item \textbf{二重盲検査読：}世界観システムを分類する研究は、科学者と当該伝統の代表者の双方が参加する審査を受けねばならない。
		
		\item \textbf{強制分析の禁止：}いかなる者もその信念についてYUCT分析を受けることを強制されてはならず、またその分析の拒否を理由に差別されてはならない。
	\end{enumerate}
	
	これらの要請は、ユネスコ「文化的多様性保護条約」（2005）およびユネスコ「神経技術倫理勧告」（2025）—精神的プライバシーと認知的自由の保護を定める—と整合する。
	
	\subsubsection{結論}
	
	YUCTは思考の構造を理解するための強力な道具を提供するが、他の道具と同様、悪用されうる。YUCTの世界観システムへの応用は、分類欲求ではなく理解欲求によって導かれるべきである。付録~$\Sigma$ に記される如く：
	
	\begin{quote}
		\itshape
		「倫理的先見性のない科学的進歩は破局へのレシピである。」
	\end{quote}
	
	ゆえに、このセクションは単なる方法論的補足ではなく、YUCTの哲学的ツールキットが世界観紛争の武器とならないことを保証する倫理的制約として機能する。
	
	% ============================================================
	% ブロック 5 終了
	% ============================================================
	
	\section{すべての科学のための統一YUCT言語}
	
	YUCTは調整効率 \(\Keff\) の概念を通じて、普遍的科学言語を創造するための根本的に新しいアプローチを提供する。これは既存の科学言語の代替ではなく、それらを単一のシステムに統合することを可能にする\textbf{メタ構造}である。
	
	\subsection{全分野のための統一メトリクス}
	
	すべての現象は三項組 \(D + I \cdot R\) とメトリクス \(\Keff\) によって記述される：
	
	\begin{table}[H]
		\centering
		\caption{異なる分野における統一YUCTメトリクスの応用}
		\label{tab:unified-metrics}
		\begin{tabular}{p{3.5cm}p{4.5cm}p{4.5cm}}
			\toprule
			\textbf{分野} & \textbf{辞書 \(D\)} & \textbf{インデックス \(I\)} \\
			\midrule
			物理学 & 自然法則、定数 & 初期条件、パラメータ \\
			生物学 & 遺伝コード、タンパク質構造 & シグナル、転写因子 \\
			社会学 & 文化的規範、制度 & 法律、政治的決断 \\
			哲学 & カテゴリー、概念 & 公理、根本原理 \\
			宗教 & 正典テキスト、儀式 & 祈り、司祭の宣言 \\
			\bottomrule
		\end{tabular}
	\end{table}
	
	\subsection{全システムのための普遍的法則}
	
	\begin{table}[H]
		\centering
		\caption{YUCTの普遍的法則}
		\label{tab:universal-laws}
		\begin{tabular}{p{5cm}p{6cm}p{5cm}}
			\toprule
			\textbf{法則} & \textbf{公式} & \textbf{応用} \\
			\midrule
			普遍誤差法則 & \(\varepsilon = \kappa_c \alpha \Keff^{-\beta}\)、\(\beta = 2/3\) & あらゆる調整システムの誤差予測 \\
			三項組 \(D + I \cdot R\) & 実在 = 辞書 + インデックス × 共鳴 & 任意のシステムを調整プロトコルとして記述 \\
			一般化ベル不等式 & \(S(\Keff) = 2 + (2\sqrt{2}-2)\cdot(1 - 2\kappa_c \alpha \Keff^{-\beta})\) & 存在論的連関性の計測 \\
			\(\Keff\) 進化法則 & \(\Keff(t) = \Keff[0] \cdot \exp(\lambda t)\) & システムの時間発展の記述 \\
			\bottomrule
		\end{tabular}
	\end{table}
	
	\subsection{統一言語の利点}
	
	伝統的科学言語には限界がある：
	\begin{itemize}
		\item 各科学分野が独自の言語を使用する。
		\item 異なる分野に共通のメトリクスがない。
		\item 異なる領域の結果を比較することが困難である。
	\end{itemize}
	
	YUCTは以下を通じてこれらの問題を解決する：
	\begin{itemize}
		\item 統一数学的装置。
		\item 全現象に共通のメトリクス。
		\item 分野横断的な結果の直接比較。
		\item 科学間の概念翻訳。
		\item 学際的理論の創出。
	\end{itemize}
	
	\subsection{統一言語の応用例}
	
	\begin{itemize}
		\item \textbf{物理学と哲学：}物理的・哲学的システムを記述するために同一メトリクスを使用。物質的および概念的システムの評価に \(\Keff\) を適用。
		\item \textbf{生物学と情報科学：}調整効率による生命システムの記述。三項組 \(D + I \cdot R\) を用いた情報システムの分析。
		\item \textbf{社会科学：}社会的システムの効率評価。調整のレンズを通した政治構造の分析。
		\item \textbf{宗教学：}事前分配辞書を持つYPSDCプロトコルとしての宗教的実践。
	\end{itemize}
	
	% ============================================================
	% ブロック 2：UCDソフトウェア実装 (短縮版)
	% ============================================================
	\subsection{ソフトウェア実装：UCD認知スペクトルメーター}
	\label{subsec:ucd-software}
	
	YUCTメトリクスの実用的応用のために、オープンソーススクリプト \texttt{run\_ucd\_telemetry.py} が開発され、リポジトリ \href{https://github.com/Alexey-Yakushev-YUCT/consciousness_meter_ucd}{consciousness\_meter\_ucd} で入手可能である。このツールはテキストストリームからの \(K_{\mathrm{eff}}\) のリアルタイム計算を実装しており、そのアーキテクチャと校正は物理学的YUCTプレプリント \cite{Yakushev2026b} に記述されている。哲学的解析の文脈では、このスクリプトは意味グラフの自動抽出と、セクション \ref{subsec:measurement-algorithm} で記述されたすべてのメトリクスの後続計算に使用される。
	
	% ============================================================
	% ブロック 2 終了
	% ============================================================
	
	\section{認識論：共鳴としての真理}
	
	\subsection{真理の定義}
	
	YUCTでは、真理は\textbf{共鳴的安定性}として定義される：
	
	\begin{definition}[真理]
		命題 \(P\) が真であるとは、それをインデックスとして受け入れることが、対応する辞書を最小誤差かつ最大予測的成功で活性化することをいう。
	\end{definition}
	
	\subsection{真理の定量的尺度}
	
	\begin{equation}
		\text{真理}(P) = \frac{R(P, D_{\text{実在}})}{\eps(P)}
	\end{equation}
	
	ここで \(R(P, D_{\text{実在}})\) は命題 \(P\) と実在辞書との間の共鳴、\(\eps(P)\) は活性化から生じる誤差である。
	
	\subsection{真理の極限定理}
	
	\begin{theorem}[真理の極限]
		完全な真理は極限 \(\Keff \to \infty\) でのみ達成可能であり、これは実在の完全な辞書に対応する。
	\end{theorem}
	
	これは\emph{いかなる}有限哲学体系も絶対的真理を主張できない理由を説明する。真理は漸近的理想である。
	
	% ============================================================
	% ブロック 3：LLMプロンプト (10.4 として挿入)
	% ============================================================
	\subsection{LLMプロンプト：哲学体系の自動計測}
	\label{sec:llm-prompt}
	
	YUCT方法論をすべての研究者が利用できるようにするため、現代の言語モデル（ChatGPT、Claude、Geminiなど）と連携するための既製プロンプトを提供する。このプロンプトはYPSDCプロトコルを実装し、AIに辞書 \(D\)（YUCT概念）とインデックス \(I\)（指示）を与え、計算されたメトリクスとして共鳴 \(R\) を出力として受け取る。
	
	\subsubsection{プロンプトテキスト}
	
	以下のテキストをコピーし、LLMとの対話に貼り付け、`[哲学書のテキストを挿入]` を分析対象テキストに置き換える。
	
	\begin{verbatim}
		あなたはヤクシェフ統一調整理論（YUCT）の専門家です。
		あなたの課題は、YPSDCプロトコルに厳密に従って哲学テキストの定量分析を行うことです。
		
		以下のテキストが与えられます：
		[哲学書のテキストを挿入]
		
		以下の手順を実行してください：
		
		1. すべての固有概念、カテゴリー、実体を抽出します。
		意味ネットワーク（概念間の接続グラフ）を構築します。
		辞書エントロピー H(D) = -sum(p_i * log2(p_i)) を計算します。ここで p_i は概念頻度です。
		
		2. システムの基本公理（通常1～5個の命題で、すべてがそこから導出される）を見つけます。
		インデックスエントロピー H(I) = -sum(q_j * log2(q_j)) を計算します。ここで q_j は公理頻度です。
		
		3. 調整効率を計算します：
		K_eff = H(D) / H(I)
		
		4. 共鳴 R = 1 /（公理から意味グラフ内の全概念への平均パス長）を推定します。
		
		5. 哲学的誤差を計算します：
		epsilon = (1/3) * 0.1 * K_eff^(-2/3)
		
		6. 一般化ベル不等式を計算します：
		S = 2 + (2*sqrt(2) - 2) * (1 - 2*(1/3)*0.1*K_eff^(-2/3))
		
		7. 以下のスケールに従ってシステムの相転移を決定します：
		- K_eff < 2：ニヒリズム／不条理
		- 2 <= K_eff < 5：懐疑主義／相対主義
		- 5 <= K_eff < 10：形而上学／体系哲学
		- 10 <= K_eff < 20：批判哲学／超越論哲学
		- K_eff >= 20：形式システム（数学）
		
		8. 誤差の主な原因（矛盾、アンチノミー、暗黙の前提）を指摘してください。
		
		9. K_eff を増加させる方法（公理の形式化、カテゴリー数の削減など）を提案してください。
		
		JSON形式で出力してください：
		{
			"text": "テキストのタイトルまたは断片",
			"H_D": 値,
			"H_I": 値,
			"K_eff": 値,
			"R": 値,
			"epsilon": 値,
			"S": 値,
			"phase_transition": "段階名",
			"main_error_sources": ["原因1", "原因2"],
			"suggestions": ["提案1", "提案2"]
		}
	\end{verbatim}
	
	\subsubsection{使用例}
	
	カント『純粋理性批判』の断片の場合、プロンプトは以下を出力する：
	
	\begin{verbatim}
		{
			"text": "純粋理性批判（断片）",
			"H_D": 8.2,
			"H_I": 0.65,
			"K_eff": 12.6,
			"R": 0.67,
			"epsilon": 0.023,
			"S": 2.81,
			"phase_transition": "批判哲学",
			"main_error_sources": ["純粋理性のアンチノミー", "統覚の統一性の暗黙の前提"],
			"suggestions": ["カテゴリーを公理系として形式化する", "先験的直観形式の数を減らす"]
		}
	\end{verbatim}
	
	\subsubsection{科学的価値}
	
	このプロンプトにより以下が可能になる：
	\begin{itemize}
		\item 哲学体系の迅速な予備比較；
		\item 隠れた矛盾の特定；
		\item 単一著者の作品における \(K_{\mathrm{eff}}\) の進化の追跡；
		\item 哲学的メトリクスデータベースの作成。
	\end{itemize}
	
	結果の検証のために、手動分析と組み合わせて使用することを推奨する。
	% ============================================================
	% ブロック 3 終了
	% ============================================================
	
	\section{倫理学：善としての \texorpdfstring{\(\Keff\)}{K\_eff} の成長}
	
	\subsection{調整倫理学}
	
	YUCTでは、倫理学は調整ダイナミクスから導出される：
	
	\begin{definition}[善]
		善とは、ネットワークの長期的・全体的調整効率 \(\Keff\) を増加させる行為である。
	\end{definition}
	
	\begin{definition}[悪]
		悪とは、\(\Keff\) を減少させるか、誤差 \(\eps\) を増加させる行為である。
	\end{definition}
	
	\subsection{倫理の定量的尺度}
	
	\begin{equation}
		\text{倫理的価値}(A) = \Delta \Keff(A) - \lambda \cdot \Delta \eps(A)
	\end{equation}
	
	ここで \(\Delta \Keff(A)\) は全体的調整効率の変化、\(\Delta \eps(A)\) は全体的誤差の変化である。
	
	\subsection{調整の定言命法}
	
	\begin{quote}
		\textit{「あなたの行為がネットワーク全体の長期的調整効率を最大化しつつ、適応を可能にするために必要な誤差レベルを維持するように行為せよ。」}
	\end{quote}
	
	% ============================================================
	% ブロック 6：教育コンポーネント (11.4 として挿入)
	% ============================================================
	\subsection{教育コンポーネント：定量的哲学の教え方}
	\label{sec:educational}
	
	YUCT方法論を教育プロセスに統合するため、以下の教材を提案する。
	
	\subsubsection{典型的課題}
	
	\begin{enumerate}[label=(\arabic*)]
		\item \textbf{プラトンのテキストの \(K_{\mathrm{eff}}\) を計測する。}主要概念と公理を抽出し、意味ネットワークを構築し、エントロピーと \(K_{\mathrm{eff}}\) を計算する。アリストテレスの結果と比較する。
		\item \textbf{哲学史における \(\eps\) の進化を分析する。}異なる時代の三つのテキストを選び、\(\eps\) を計算し、誤差がなぜ変化したかを説明する。
		\item \textbf{哲学的危機の診断。}ポストモダンのテキスト断片から \(K_{\mathrm{eff}}\) を計算し、システムがニヒリズムの段階にあるかどうかを判定する。
	\end{enumerate}
	
	\subsubsection{実験演習の例}
	
	\textbf{テーマ：}「形而上学的体系の調整効率の比較分析」。
	
	\begin{enumerate}
		\item 学生は三つの断片（デカルト、カント、ニーチェ）を得る。
		\item 意味解析（または手動）を用いて辞書とインデックスを構築する。
		\item \(K_{\mathrm{eff}}\)、\(\eps\)、\(R\)、\(S\) を計算する。
		\item グラフを作成し、各システムの深さと連関性について結論を導く。
		\item 各システムの \(K_{\mathrm{eff}}\) をどのように増加させられるかを議論する。
	\end{enumerate}
	
	\subsubsection{指導者向け方法論的推奨}
	
	\begin{itemize}
		\item 簡単なテキストから始める（例えばデカルトの『方法序説』）。
		\item 公理の同定にはグループ討論を利用する。
		\item 学生が独自のLLMプロンプトを開発するよう奨励する。
		\item 検証のために手動計算と自動計算を比較する。
	\end{itemize}
	% ============================================================
	% ブロック 6 終了
	% ============================================================
	
	\section{美学：圧縮としての美}
	
	\subsection{美の定義}
	
	YUCTでは、美は\textbf{最大圧縮のインデックス}として定義される：
	
	\begin{definition}[美]
		美とは、準備された辞書 \(D\) に出会ったとき、最小エネルギー消費で最大振幅の共鳴 \(R\) を引き起こす極度の圧縮のインデックス \(I\) である。
	\end{definition}
	
	\subsection{美の定量的尺度}
	
	\begin{equation}
		\text{美}(\mathcal{A}) = \frac{\Delta \Keff(\text{観衆}) \cdot R(D_{\text{観}}, I_{\text{美}})}{L(I_{\text{美}})}
	\end{equation}
	
	ここで \(\Delta \Keff\) は調整効率の増加、\(R\) は共鳴的適合性、\(L(I)\) はインデックスの長さである。
	
	\subsection{美的三項組}
	
	\begin{itemize}
		\item \textbf{醜}—共鳴しないか辞書を破壊するインデックス。
		\item \textbf{美}—既存の辞書と共鳴し、圧縮の突然の増加をもたらすインデックス。
		\item \textbf{崇高}—情報密度が辞書の現在の容量を超え、その拡張を強いるインデックス。
	\end{itemize}
	
	\section{政治哲学：調整と権力}
	
	\subsection{調整システムとしての国家}
	
	YUCTでは、国家は以下の調整システムである：
	
	\begin{itemize}
		\item \textbf{辞書} \(D_{\text{国家}}\)—法律、制度、規範。
		\item \textbf{インデックス} \(I_{\text{国家}}\)—政治的決定、法律、勅令。
		\item \textbf{共鳴} \(R_{\text{国家}}\)—正当性、統治効率。
	\end{itemize}
	
	\subsection{凍結辞書としての専制}
	
	専制は、単一の硬直した辞書 \(D_{\text{専制}}\) を押し付け、あらゆる局所的更新を禁止する試みである。これは高い瞬間的 \(\Keff\) を与えるが、三つの構造的弱点をもたらす：
	
	\begin{enumerate}
		\item \textbf{辞書の凍結：}システムが新しいインデックスに適応できない。
		\item \textbf{隠れた誤差の蓄積：} \(\eps\) が成長し、抑圧により多くのエネルギーを要する。
		\item \textbf{エネルギーの枯渇：}行動予算全体が自己維持に費やされる。
	\end{enumerate}
	
	\subsection{分散誤差としての自由}
	
	自由は、システムがグローバル辞書 \(D_{\text{グローバル}}\) から逸脱する局所辞書 \(D_i\) を維持する能力である。これは贅沢ではなく、\textbf{生存メカニズム}である。辞書の多様性は潜在的な適応の分散リザーバーとして機能する。
	
	\section{歴史的先駆者：調整の直観}
	
	\subsection{ライプニッツ：モナドとしての辞書}
	
	ゴットフリート・ヴィルヘルム・ライプニッツは『モナドロギー』（1714）において、宇宙をモナド—閉じた「窓のない」実体で、信号を交換しないが\textbf{予定調和}によって完全に同期している—から成ると記述した。
	
	YUCTでは、これは哲学的形式におけるYPSDCプロトコルである。モナドは内部辞書 \(D\)、現在状態 \(I\)、およびグローバル秩序との共鳴 \(R\) を持つ調整ノードである。予定調和はオフライン段階であり、辞書は創造の瞬間にすべてのモナドに分配される。
	
	\subsection{EPRパラドックス：オフライン辞書の必要性の証明}
	
	1935年、アインシュタイン、ポドルスキー、ローゼンは量子力学が「幽霊的な遠隔作用」を暗示することを示した。
	
	YUCTはこのパラドックスを解決する：もつれ合った粒子間の相関はシグナルではなく、\emph{もつれの瞬間にオフラインで分配された共通辞書の活性化}である。
	
	\subsection{シャノン：情報と意味の分離}
	
	クロード・シャノンは1948年に情報を意味から分離した。YUCTはさらに進む：\emph{調整}を\emph{情報}から分離する。情報はインデックスであり、調整は辞書活性化のプロセスである。
	
	% ============================================================
	% 新章：同型構造 (15 として挿入)
	% ============================================================
	\section{科学間の架け橋としての同型構造}
	\label{sec:isomorphisms}
	
	調整アプローチの最も強力な結果の一つは、異なる科学分野における\textbf{同型構造}の発見である。ここでの同型は単なる類推ではなく、根本的に異なるシステムを記述する関数依存関係の\emph{定量的一致}を意味する。
	
	\subsection{調整同型の定義}
	
	\begin{definition}[調整同型]
		二つのシステム \(A\) と \(B\) が調整同型であるとは、以下の全単射写像が存在することをいう：
		\[
		\Phi: (D_A, I_A, R_A, \Keff[A], \eps_A) \mapsto (D_B, I_B, R_B, \Keff[B], \eps_B),
		\]
		これが以下を保存する：
		\begin{enumerate}
			\item 三項組 \(D+I\cdot R\) の構造、
			\item 普遍誤差法則 \(\eps = \kappac \alpha \Keff^{-\betaf}\)（\(\betaf = 2/3\)）、
			\item 臨界閾値 \(\Rcrit\) と相転移。
		\end{enumerate}
	\end{definition}
	
	\subsection{同型構造の例}
	
	表 \ref{tab:isomorphisms} は、異なる科学分野からのいくつかの顕著な例を示し、同一の調整ダイナミクスを実証する。
	
	\begin{table}[H]
		\centering
		\caption{異なる科学分野における同型構造}
		\label{tab:isomorphisms}
		\footnotesize
		\begin{tabular}{p{2.8cm}p{3cm}p{3cm}p{3.5cm}}
			\toprule
			\textbf{分野} & \textbf{辞書 \(D\)} & \textbf{インデックス \(I\)} & \textbf{誤差法則} \\
			\midrule
			物理学（量子重力） & 状態多様体 \(\mathcal{M}\) & ラグランジアン \(\mathcal{L}\) & \(\delta g_{\mu\nu} \propto \Keff^{-2/3}\) \\
			遺伝学 & 遺伝コード & 転写シグナル & \(\mu \propto K_{\mathrm{repair}}^{-2/3}\) \\
			経済学 & 金融商品 & 市場指数 & \(\sigma_{\mathrm{GDP}} \propto W^{-2/3}\) \\
			社会学 & 文化的規範 & 政治的決断 & \(\Delta \text{調整} \propto N^{-2/3}\) \\
			哲学 & 概念辞書 & 基本公理 & \(\eps_{\text{哲}} = \kappac \alpha \Keff^{-2/3}\) \\
			\bottomrule
		\end{tabular}
	\end{table}
	
	\subsection{科学分野間の定量的関係}
	
	同型により、異なる科学のパラメータ間に\textbf{正確な定量的関係}を確立できる。例えば、遺伝学における突然変異率 \(\mu\) と経済学における経済変動性 \(\sigma_{\mathrm{GDP}}\) は同じべき法則に従う：
	
	\[
	\mu \propto K_{\mathrm{repair}}^{-2/3}, \qquad \sigma_{\mathrm{GDP}} \propto W^{-2/3}.
	\]
	
	\(K_{\mathrm{repair}}\) をDNA修復効率で表し、\(W\) を取引量で表すと、それらの相対的変化が普遍定数 \(\betaf = 2/3\) によって結ばれていることが分かる。これは経験的偶然ではなく、統一調整的性質の結果である。
	
	\subsection{調整システムとしての哲学}
	
	同型に照らして、哲学は孤立した学問でなくなる。それは以下のパラメータを持つ調整システムとなる：
	
	\begin{itemize}
		\item \textbf{辞書} \(D_{\text{哲}}\)—概念、カテゴリー、存在論的実体の集合。
		\item \textbf{インデックス} \(I_{\text{哲}}\)—基本公理（例えばデカルトの \textit{我思う、ゆえに我あり}）。
		\item \textbf{共鳴} \(R_{\text{哲}}\)—システムの連関性、一公理が全体の存在論を決定する程度。
		\item \textbf{調整効率} \(\Keff[\text{哲}] = H(D)/H(I)\)。
		\item \textbf{誤差} \(\eps_{\text{哲}} = \kappac \alpha_{\text{哲}} \Keff[\text{哲}]^{-2/3}\)。
	\end{itemize}
	
	このシステムは遺伝的、経済的、物理的システムと同じ法則に従う。したがって、他の科学で開発された方法がそれに適用可能である。
	
	\subsection{相転移による真理の達成}
	
	YUCTでは、真理は絶対状態ではなく、極限 \(\Keff \to \infty\) で達成され、誤差 \(\eps \to 0\) となる。しかし有限システムでは、真理は\textbf{相転移}—臨界閾値 \(\Rcrit\) を通じた調整効率のジャンプ—に対応する。
	
	哲学にとってこれは、\emph{真理は知識の蓄積ではなく、辞書の質的再構築によって達成される}ことを意味する。その際 \(\Keff[\text{哲}]\) が閾値を超える。このような転移は以下の場合に可能である：
	
	\begin{enumerate}
		\item \textbf{インデックス} \(I\) が十分に短く強力になる（一公理が全システムを生成）。
		\item \textbf{辞書} \(D\) が十分に連関的で無矛盾になる。
		\item \textbf{共鳴} \(R\) が最大に達する（各概念が全システムを活性化）。
	\end{enumerate}
	
	同型構造の分析は、そのような転移が哲学史において既に起こったことを示す：プラトン総合、デカルト転回、カント革命。YUCT用語では、それぞれが \(\Keff[\text{哲}]\) の一桁以上のジャンプである。
	
	\subsection{同型が哲学の洗練にどう役立つか}
	
	同型の知識により以下が可能になる：
	
	\begin{enumerate}
		\item \textbf{方法の移転：}遺伝学における \(\Keff\) 増加のために開発された方法（例：コドン辞書最適化）や経済学（インデックスエントロピーの削減）を哲学システムに適用できる。
		\item \textbf{危機の診断：} \(\eps_{\text{哲}}\) の値から、哲学システムが安定状態にあるか崩壊状態（ポストモダン）にあるかを判定できる。
		\item \textbf{発展の予測：} \(\Keff[\text{哲}]\) のダイナミクスを知れば、システムがいつ閾値に達し相転移を起こすかを予測できる。
		\item \textbf{標的型介入：}例えば、新公理の導入や概念装置の再構築が \(\Keff[\text{哲}]\) を高め、新たな総合をもたらす可能性がある。
	\end{enumerate}
	
	\subsection{例：デカルトからYUCTへ}
	
	デカルトから現代YUCTへの哲学的調整の進化を考察する。
	
	\begin{itemize}
		\item \textbf{デカルト：} \(I = \{\textit{我思う、ゆえに我あり}\}\)、\(D\)—デカルト形而上学、\(\Keff \approx 12\)、\(\eps \approx 0.028\)。
		\item \textbf{カント：} \(I = \{\text{先験的直観形式、悟性カテゴリー}\}\)、\(D\)—超越論哲学、\(\Keff \approx 10\)、\(\eps \approx 0.033\)。
		\item \textbf{ポストモダン：} \(I\) 欠如、\(D\) 断片化、\(\Keff \approx 2\)、\(\eps \approx 0.10\)。
		\item \textbf{YUCT：} \(I = \{\text{調整としての存在論的原初概念}\}\)、\(D\)—統一調整辞書、\(\Keff \gg 10\)、\(\eps \to 0\)。
	\end{itemize}
	
	この系列は、哲学が \(\eps\) が最小で \(\Keff\) が最大の状態に到達できることを示す。同型分析により、過去に \(\Keff\) のジャンプをもたらした辞書とインデックスの具体的な変更を理解し、その知識をさらなる洗練に応用できる。
	
	\subsection{統一の実用的利益}
	
	調整同型による科学の統一は以下をもたらす：
	
	\begin{enumerate}
		\item \textbf{統一言語：}すべての科学が一つの調整言語を話し、学際的障壁を取り除く。
		\item \textbf{相互豊化：}物理学の方法が哲学に、生物学の方法が経済学に応用可能。
		\item \textbf{予測力：}一つのシステムの誤差法則を知れば、同型システムの挙動が予測できる。
		\item \textbf{進歩の加速：}各分野で法則を再発見する代わりに、既知の構造を移転する。
		\item \textbf{真理の到達可能性：}哲学は他の調整システムと同様、\(\Keff\) の意図的増加を通じて最小誤差状態に到達できる。
	\end{enumerate}
	
	\begin{quote}
		\itshape
		「同型は単なる奇妙な偶然ではない。それはすべての科学分野が同一の調整現実の異なる投影であることの証明である。それらを統合することで、多様性を失うのではなく、深みを得るのである。」
	\end{quote}
	
	% ============================================================
	% 新章終了
	% ============================================================
	
	% ============================================================
	% ブロック 7：批判的分析 (セクション16)
	% ============================================================
	\section{批判的分析：適用限界と可能な反論}
	\label{sec:critical-analysis}
	
	いかなる新方法論もその限界を明確に定義しなければならない。以下に主な限界と可能な反論およびそれらへの応答を示す。
	
	\subsection{方法論の限界}
	
	\begin{enumerate}
		\item \textbf{言語と翻訳への依存。}意味解析は選択された言語に敏感である。翻訳は概念頻度や接続構造を歪める可能性がある。\textit{解決策：}原文を使用するか、複数言語で分析を実施する。
		\item \textbf{公理同定の主観性。}研究者によって異なる公理を同定する可能性がある。\textit{解決策：}形式基準（最小長さ、最大生成能力）を使用し、専門家間検証を実施する。
		\item \textbf{暗黙の前提。}多くの哲学体系はテキストに表現されていない暗黙の前提を含む。\textit{解決策：}歴史・哲学的コンテキストで分析を補完する。
		\item \textbf{テキスト量。}統計的有意性には十分なボリュームが必要である。短い断片は信頼できない推定値を与える。\textit{解決策：}コーパス全体を分析する。
	\end{enumerate}
	
	\subsection{潜在的な誤差源}
	
	\begin{itemize}
		\item 解析エラー（概念の誤認識）。
		\item 意味的接続におけるノイズ（偽の相関）。
		\item 著者のスタイルが頻度特性に与える影響。
		\item ジャンルの違い（論文 vs. 対話 vs. 箴言）。
	\end{itemize}
	
	\subsection{代替アプローチとの比較}
	
	\begin{itemize}
		\item \textbf{形式論理}は統語的厳密性を提供するが、意味論とシステムダイナミクスをカバーしない。
		\item \textbf{解釈学}は意味を深く分析するが、定量的尺度を提供しない。
		\item \textbf{議論理論}は議論の強さを計測するが、システム全体の連関性は計測しない。
	\end{itemize}
	YUCTは妥協を提供する：意味ネットワークを通じて意味的深さを維持しつつ、定量的メトリクスを提供する。
	
	\subsection{潜在的反論への応答}
	
	\begin{quote}
		\textit{「哲学は計測できない。物理学ではない。」}
	\end{quote}
	— 構造を持ち記述可能なものはすべて計測可能である。YUCTは哲学を還元するのではなく、その内部組織を形式化する。
	
	\begin{quote}
		\textit{「これは還元主義だ。哲学が数字に還元される。」}
	\end{quote}
	— これは還元ではなく補完である。数字は新しい視点を提供するが、内容分析を無効にしない。
	
	\begin{quote}
		\textit{「これらのメトリクスが意味を持つという証拠はどこにある？」}
	\end{quote}
	— 普遍誤差法則 \(\eps \propto K_{\mathrm{eff}}^{-2/3}\) は物理学、生物学、社会学において50オーダー以上にわたって確認されている。それを哲学に拡張することは自然な一般化である。
	% ============================================================
	% ブロック 7 終了
	% ============================================================
	
	% ============================================================
	% ブロック 8：学際的連携 (セクション17)
	% ============================================================
	\section{学際的連携：物理学から認知科学へ}
	\label{sec:interdisciplinary}
	
	YUCTはメトリクスを分野間で移転可能にする統一言語を提供する。これにより学際的研究に新たな可能性が開かれる。
	
	\subsection{認知科学における哲学的メトリクスの応用}
	
	認知科学では、\(K_{\mathrm{eff}}\) は認知課題の複雑性評価に、\(\eps\) は認知的バイアスの計測に使用できる。意識は \(K_{\mathrm{eff}} > K_{\mathrm{crit}} \approx 8.5\) を持つシステムの状態として解釈される（定理6参照）。意識の「ハード・プロブレム」（チャルマーズ）は調整的解消を得る：主観的経験は内部辞書間の高共鳴接続性の効果として生じる。
	
	\subsection{情報理論との連携}
	
	シャノン理論は、共鳴 \(R = 1\)（増幅なし）の場合のYUCTの特殊例である。YUCTは調整効率と共鳴の概念を追加することでシャノンを一般化する。これにより情報の量だけでなく深さも記述できる。
	
	\subsection{量子力学との連携}
	
	一般化ベル不等式 \(S(\Keff)\) は、量子非局所性が調整連関性の特殊例であることを示す。\(K_{\mathrm{eff}} \to \infty\) のとき \(S = 2\sqrt{2}\) を得、これは最大エンタングルメントに対応する。したがってYUCTは存在論、認識論、物理学を結びつける。
	
	\subsection{言語学との連携}
	
	意味ネットワークと辞書エントロピーは言語的世界観の比較に応用できる。言語間の \(K_{\mathrm{eff}}\) の違いは文化的思考特性を説明するかもしれない。
	
	\subsection{美術史との連携}
	
	作品の美学的価値は共鳴 \(R\) と観衆における \(\Delta K_{\mathrm{eff}}\) の増加を通じて評価できる。これは定量的美学への道を開く。
	% ============================================================
	% ブロック 8 終了
	% ============================================================
	
	% ============================================================
	% ブロック 9：将来展望 (セクション18)
	% ============================================================
	\section{方法論発展の将来展望}
	\label{sec:future-directions}
	
	提案されたアプローチは、さらなる研究と実用的応用に広い可能性を開く。
	
	\subsection{方法論拡張の計画}
	
	\begin{enumerate}
		\item \textbf{分析の自動化。}哲学的テキストの大量意味解析のためのソフトウェアスイートの開発。
		\item \textbf{メトリクスデータベース。}すべての重要な哲学的著作の \(K_{\mathrm{eff}}\)、\(\eps\)、\(R\)、\(S\) のオープンデータベースの構築。
		\item \textbf{可視化。}哲学的思考の進化のインタラクティブマップの開発。
		\item \textbf{LLMとの統合。}哲学的解析のための専門AIアシスタントの作成。
	\end{enumerate}
	
	\subsection{さらなる研究の可能な方向性}
	
	\begin{itemize}
		\item テキストが他の言語に翻訳されると \(K_{\mathrm{eff}}\) はどう変化するか？
		\item \(K_{\mathrm{eff}}\) と美学的価値の間に相関はあるか？
		\item \(\Keff\) のダイナミクスから新しい哲学体系の出現を予測できるか？
		\item 文化的コンテキストは \(\alpha_{\text{哲}}\) にどのように影響するか？
		\item 方法論は非西洋哲学の伝統に適用可能か？
	\end{itemize}
	
	\subsection{期待される成果}
	
	\textbf{短期（1～2年）：}
	\begin{itemize}
		\item オープンソースアナライザの実用的プロトタイプの作成。
		\item 50の主要テキストの \(K_{\mathrm{eff}}\) 初期データベース。
		\item いくつかの教育ケーススタディの公開。
	\end{itemize}
	
	\textbf{長期（3～5年）：}
	\begin{itemize}
		\item 哲学の完全な実験科学への転換。
		\item 調整哲学の国際研究ネットワークの創設。
		\item YUCT方法論の大学カリキュラムへの組み込み。
	\end{itemize}
	% ============================================================
	% ブロック 9 終了
	% ============================================================
	
	% ============================================================
	% ブロック 10：用語集 (セクション19)
	% ============================================================
	\section{用語集と専門用語}
	\label{sec:glossary}
	
	読者の便宜のため、本稿で使用されるすべての主要メトリクスと概念の正確な定義を提供する。
	
	\begin{longtable}{p{3.5cm}p{5cm}p{4cm}}
		\toprule
		\textbf{用語} & \textbf{記号} & \textbf{定義} \\
		\midrule
		\endfirsthead
		\toprule
		\textbf{用語} & \textbf{記号} & \textbf{定義} \\
		\midrule
		\endhead
		辞書（哲学的） & \(D\) & システムが活性化できるすべての概念、カテゴリー、判断の集合。 \\
		インデックス & \(I\) & 基本公理、システム全体がそこから展開される。 \\
		共鳴 & \(R\) & インデックスが辞書を活性化する完全性の尺度；意味ネットワークでは平均パス長の逆数。 \\
		調整効率 & \(\Keff\) & 辞書エントロピーとインデックスエントロピーの比：\(K_{\mathrm{eff}} = H(D)/H(I)\)。 \\
		哲学的誤差 & \(\eps\) & 内部矛盾の尺度：\(\eps = \kappac \cdot \alpha \cdot K_{\mathrm{eff}}^{-2/3}\)。 \\
		一般化ベル不等式 & \(S\) & 存在論的連関性の尺度：\(S = 2 + (2\sqrt{2}-2)(1 - 2\kappac\alpha K_{\mathrm{eff}}^{-2/3})\)。 \\
		個人調整行列 & \(\Cpers\) & 人格（「魂」）を定義する結合定数と共鳴因子の個別集合。 \\
		辞書エントロピー & \(H(D)\) & \( -\sum p_i \log_2 p_i\)、ここで \(p_i\) は概念頻度。 \\
		インデックスエントロピー & \(H(I)\) & \( -\sum q_j \log_2 q_j\)、ここで \(q_j\) は公理頻度。 \\
		存在論スペクトル & \(\alpha_D\) & 単位時間あたりの語彙的豊かさ（UCDで使用）。 \\
		情報スペクトル & \(\beta_I\) & 認知的処理の深さ（UCD）。 \\
		共鳴スペクトル & \(\gamma_R\) & 自然なリズム変動性（UCD）。 \\
		意識閾値 & \(K_{\mathrm{crit}}\) & \(K_{\mathrm{eff}} = 8.5\) の臨界値、これを超えるとシステムは自己認識を得る。 \\
		\bottomrule
		\caption{YUCT哲学的解析の主要用語集。}
		\label{tab:glossary}
	\end{longtable}
	
	\subsection{伝統的哲学概念とYUCTメトリクスの対応表}
	
	\begin{table}[H]
		\centering
		\caption{哲学概念とYUCTメトリクスの対応}
		\label{tab:philosophy-metrics}
		\begin{tabular}{p{3.5cm}p{3.5cm}p{4cm}}
			\toprule
			\textbf{伝統的概念} & \textbf{YUCTメトリクス} & \textbf{説明} \\
			\midrule
			システムの深さ & \(\log(K_{\mathrm{eff}})\) & 調整効率の対数 \\
			矛盾性 & \(\eps\) & 哲学的誤差 \\
			連関性 & \(S\) & 一般化ベル不等式 \\
			魂 & \(\Cpers\) & 個人調整行列 \\
			真理 & 共鳴的安定性 & \(R(P, D_{\text{実在}}) / \eps(P)\) \\
			善 & \(\Delta \Keff\) & 全体的調整効率の増加 \\
			美 & 圧縮 & \( \Delta \Keff(\text{観衆}) \cdot R / L(I) \) \\
			自由 & \(\partial \Cpers / \partial I\) & インデックス選択による行列変更能力 \\
			\bottomrule
		\end{tabular}
	\end{table}
	% ============================================================
	% ブロック 10 終了
	% ============================================================
	
	% ============================================================
	% ブロック 11：可視化 (セクション20)
	% ============================================================
	\section{メトリクスの可視化例}
	\label{sec:visualizations}
	
	分析結果の視覚的提示のために、以下の可視化タイプを提案する。
	
	\subsection{哲学史における \(K_{\mathrm{eff}}\) の進化グラフ}
	
	\begin{figure}[H]
		\centering
		\begin{tikzpicture}
			\begin{axis}[
				width=0.9\textwidth,
				height=0.5\textwidth,
				xlabel={時間（時代）},
				ylabel={\(K_{\mathrm{eff}}\)},
				xtick={1,2,3,4,5,6,7},
				xticklabels={古代, 中世, ルネサンス, 近代, 19世紀, 20世紀, 21世紀},
				ymin=0, ymax=25,
				grid=both,
				legend pos=north west,
				]
				\addplot[thick, blue, mark=*] coordinates {
					(1,5) (2,4) (3,6) (4,10) (5,15) (6,8) (7,20)
				};
				\addlegendentry{\(K_{\mathrm{eff}}\)}
				\draw[dashed, red] (axis cs:1,8.5) -- (axis cs:7,8.5);
				\node at (axis cs:4,9) {\(K_{\mathrm{crit}} = 8.5\)};
			\end{axis}
		\end{tikzpicture}
		\caption{哲学史における調整効率の進化。19世紀の上昇と20世紀（ポストモダン）の下降が可視。}
		\label{fig:keff-evolution}
	\end{figure}
	
	\subsection{哲学的誤差 \(\eps\) の分布図}
	
	\begin{figure}[H]
		\centering
		\begin{tikzpicture}
			\begin{axis}[
				width=0.8\textwidth,
				height=0.4\textwidth,
				xlabel={\(\eps\)},
				ylabel={頻度},
				ymin=0,
				grid=both,
				]
				\addplot[hist={bins=10}, fill=blue!30] coordinates {
					(0.01,0) (0.02,5) (0.03,12) (0.04,15) (0.05,8) (0.06,4) (0.08,2) (0.10,1)
				};
			\end{axis}
		\end{tikzpicture}
		\caption{50の主要哲学テキストの \(\eps\) 分布。ピークは0.03–0.04で、体系的な哲学体系に対応。}
		\label{fig:epsilon-distribution}
	\end{figure}
	
	\subsection{パラメータ間の関係スキーマ}
	
	\begin{figure}[H]
		\centering
		\begin{tikzpicture}[
			node distance=4cm and 3.5cm,
			auto,
			every node/.style={
				rectangle,
				minimum width=3.2cm,
				minimum height=1.0cm,
				align=center,
				font=\small\bfseries,
				rounded corners=3pt,
				line width=0.8pt
			},
			every path/.style={thick, -stealth, line width=1.2pt}
			]
			% 上段：D, I, R – draw を明示追加
			\node[fill=blue!15, draw] (D) {辞書 \(D\)};
			\node[fill=green!15, draw, below left of=D, xshift=-3.8cm, yshift=-1.2cm] (I) {インデックス \(I\)};
			\node[fill=orange!15, draw, below right of=D, xshift=2.5cm, yshift=-1.2cm] (R) {共鳴 \(R\)};
			
			% 中段：K_eff
			\node[fill=red!15, draw, below of=D, yshift=-2.5cm] (K) {\(K_{\mathrm{eff}} = \dfrac{H(D)}{H(I)}\)};
			
			% 下段：eps と S
			\node[fill=purple!15, draw, below left of=K, xshift=-3cm, yshift=-2.5cm] (eps) {\(\eps = \kappa_c \alpha K^{-2/3}\)};
			\node[fill=teal!15, draw, below right of=K, xshift=3cm, yshift=-2.5cm] (S) {\(S = f(K)\)};
			
			% 矢印とラベル – 枠なし
			\draw[->] (D) -- node[above, sloped, font=\small] {エントロピー} (K);
			\draw[->] (I) -- node[above, sloped, font=\small] {エントロピー} (K);
			\draw[->] (D) -- node[above, sloped, font=\small] {構造} (R);
			\draw[->] (I) -- node[above, sloped, font=\small] {活性化} (R);
			\draw[->] (K) -- node[left, font=\small] {誤差} (eps);
			\draw[->] (K) -- node[above, sloped, font=\small] {連関性} (S);
			\draw[->] (R) -- node[above, sloped, font=\small] {増幅} (S);
			% 追加接続：R から K_eff への影響（任意）
			\draw[->, dashed, gray!60] (R) -- node[right, font=\small] {共鳴} (K);
		\end{tikzpicture}
		\caption{主要YUCTメトリクス間の関係スキーマ。直接および媒介接続を示す：辞書 \(D\) とインデックス \(I\) が \(K_{\mathrm{eff}}\) を決定し、それが誤差 \(\eps\) と存在論的連関性 \(S\) に影響；共鳴 \(R\) は連関性を増幅し調整効率に影響を与える。}
		\label{fig:parameter-relations}
	\end{figure}
	
	% ============================================================
	% 新サブセクション：相関行列 S と自動ベル非局所性チェック
	% ============================================================
	\subsection{相関行列 \(S\) と自動ベル非局所性チェック}
	\label{subsec:bell-matrix}
	
	複数の哲学システム（または単一システムの断片）間の存在論的連関性の定量的評価のために、一般化ベル不等式に基づく\textbf{相関行列 \(S_{ij}\)} を構築する。
	
	\subsubsection{行列の形式定義}
	
	システムの集合 \(\{\mathcal{S}_1, \mathcal{S}_2, \dots, \mathcal{S}_k\}\) が与えられたとき、各ペア \((\mathcal{S}_i, \mathcal{S}_j)\) についてパラメータ \(S_{ij}\) を以下のように計算する：
	
	\[
	S_{ij} = 2 + (2\sqrt{2} - 2) \cdot \left(1 - 2\kappac \alpha \, \Keff[ij]^{-2/3}\right),
	\]
	
	ここで \(\Keff[ij]\) はシステム \(i\) と \(j\) 間の調整効率であり、以下で定義される：
	
	\[
	\Keff[ij] = \frac{H(D_i \cap D_j)}{H(I_i \cap I_j)},
	\]
	
	すなわち、辞書の交差のエントロピーとインデックスの交差のエントロピーの比である。システムが共通の概念や公理を持たない場合、\(\Keff[ij] = 1\) かつ \(S_{ij} = 2\)（古典的極限）となる。最大値 \(S_{ij} = 2\sqrt{2} \approx 2.828\) は理想的な連関性で達成される。
	
	\subsubsection{行列の自動生成}
	
	行列 \(S\) 構築のアルゴリズム：
	
	\begin{enumerate}
		\item 各テキスト（またはセクション）について、意味グラフを構築し、公理コアを抽出する（セクション \ref{sec:auto-basis} の方法による）。
		\item 各ペアについて、辞書とインデックスの交差、そのエントロピー、次に \(\Keff[ij]\) と \(S_{ij}\) を計算する。
		\item 結果をヒートマップとして表示し、セルの色が \(S_{ij}\) の値（2.0 から 2.828）に対応し、対角線には \(S_{ii} = 2\sqrt{2}\)（最大自己連関性）が配置される。
	\end{enumerate}
	
	四人の哲学者の例行列：
	
	\[
	\begin{array}{c|cccc}
		& \text{プラトン} & \text{デカルト} & \text{カント} & \text{ニーチェ} \\
		\hline
		\text{プラトン} & 2.828 & 2.61 & 2.45 & 2.18 \\
		\text{デカルト} & 2.61 & 2.828 & 2.58 & 2.22 \\
		\text{カント}   & 2.45 & 2.58 & 2.828 & 2.35 \\
		\text{ニーチェ}  & 2.18 & 2.22 & 2.35 & 2.828 \\
	\end{array}
	\]
	
	2.5 以上の値は強い存在論的連関性（共通のカテゴリー、類似の公理）を示す。2.3 以下の値は辞書の有意な乖離を示す。
	
	\subsubsection{非局所性テスト}
	
	行列 \(S\) は哲学的言説の非局所性仮説の自動テストを可能にする。あるシステムのペアについて \(S_{ij} > 2\) が統計的に有意である場合、それらの結論は独立性では説明できず、共通の辞書の存在—哲学における量子もつれの類似物—を必要とすることを意味する。
	
	\subsubsection{ソフトウェア実装}
	
	以下はテキストの集合から \(S\) 行列を計算するPythonコードである。自己完結のために、文から意味グラフを構築する簡単な関数も含まれている。
	
	\begin{lstlisting}[caption={相関行列 S 構築関数（可視化付き）}, label={lst:bell-matrix}]
		import numpy as np
		import networkx as nx
		import seaborn as sns
		import matplotlib.pyplot as plt
		
		def build_semantic_graph(text_corpus):
		"""
		意味グラフ構築の簡易実装。
		text_corpus: 文のリスト（各文は単語のリスト）。
		同じ文に共起する単語を辺で結ぶ無向グラフを返す（各文内で完全連結）。
		"""
		G = nx.Graph()
		for sent in text_corpus:
		unique_words = list(set(sent))
		for i, w1 in enumerate(unique_words):
		for w2 in unique_words[i+1:]:
		if w1 != w2:
		G.add_edge(w1, w2)
		return G
		
		def compute_S_matrix(texts, M_axioms=5):
		"""
		texts: コーパスのリスト（各コーパスは文のリスト）。
		行列 S (numpy array) を返す。
		"""
		n = len(texts)
		S_mat = np.zeros((n, n))
		# 事前にグラフと公理を構築
		graphs = []
		dicts = []
		idxs = []
		for txt in texts:
		G = build_semantic_graph(txt)
		cent = nx.eigenvector_centrality_numpy(G)
		axioms = extract_independent_basis_sentences(txt, cent, M_axioms)
		graphs.append(G)
		dicts.append(set(G.nodes()))
		idxs.append(set(word for ax in axioms for word in ax if word in G))
		for i in range(n):
		for j in range(n):
		if i == j:
		S_mat[i,j] = 2 * np.sqrt(2)
		continue
		common_D = dicts[i].intersection(dicts[j])
		common_I = idxs[i].intersection(idxs[j])
		if not common_D or not common_I:
		S_mat[i,j] = 2.0
		continue
		# 交差のエントロピー
		# 重みにはグラフ i の固有ベクトル中心性を使用
		cent = nx.eigenvector_centrality_numpy(graphs[i])
		p = np.array([cent.get(w, 0.0) for w in common_D])
		p = p / (p.sum() + 1e-12)
		H_D = -np.sum(p * np.log2(p + 1e-12))
		# 注意: インデックス重みには一様分布を使用。
		# これは K_ij の保守的推定を与える。より高い精度のためには、
		# 各交差公理集合に対して compute_axiom_coverages による
		# q_j の計算を推奨する。
		q = np.ones(len(common_I)) / len(common_I)
		H_I = -np.sum(q * np.log2(q + 1e-12))
		K_ij = H_D / H_I if H_I > 0 else 1.0
		S_mat[i,j] = 2 + (2*np.sqrt(2)-2)*(1 - 2/3 * 0.1 * K_ij**(-2/3))
		return S_mat
		
		# ヒートマップ描画の例
		def plot_S_matrix(S_mat, labels):
		sns.heatmap(S_mat, annot=True, fmt='.2f', xticklabels=labels, yticklabels=labels,
		cmap='coolwarm', vmin=2.0, vmax=2.828)
		plt.title('存在論的連関性行列 S')
		plt.show()
	\end{lstlisting}
	
	この行列の使用により、哲学システムの非局所性チェックが完全に自動化され、再現可能になる。
	
	% ============================================================
	% ブロック 11 終了
	% ============================================================
	
	\section{結論：精密科学としての哲学}
	
	我々は以下を示した：
	
	\begin{enumerate}
		\item \textbf{哲学は計測可能である。}いかなる哲学体系も \(\Keff\)、\(\eps\)、\(R\)、\(S\) によって定量的に記述できる。
		\item \textbf{哲学的袋小路は診断可能である。}三項組 \(D + I \cdot R\) のいずれかの要素の欠如はシステム崩壊をもたらす。
		\item \textbf{哲学史は \(\Keff\) の進化である。}神話から形而上学へ、形而上学から数学へ。
		\item \textbf{倫理学・美学・政治学は調整から導出される。}善、美、自由は定量的定義を持つ。
		\item \textbf{魂は厳密な定義を得る。}魂は個別調整行列 \(\Cpers\) である。
		\item \textbf{YUCTは「もう一つの理論」ではない。}それは、それを定式化する理論を含む、あらゆる理論の基礎にある\textbf{プロトコル}である。
	\end{enumerate}
	
	提案されたアプローチにより、哲学を計測可能な学問として捉えることが可能となり、哲学体系の比較分析、危機の診断、発展の予測に新たな機会が開かれる。
	
	\begin{center}
		\fbox{\parbox{0.85\textwidth}{\centering\Large\bfseries
				哲学はもはや推論の芸術ではない。\\
				それは調整に関する厳密な科学となる。\\
				歴史に初めて—哲学は計測可能となる。
		}}
	\end{center}
	
	\subsection*{謝辞}
	
	著者は、調整、予定調和、非局所性についての直観がYUCTにおいて完成を見たライプニッツ、シャノン、アインシュタイン、ポドルスキー、ローゼンに感謝する。
	
	\subsection*{データ利用可能性}
	
	すべての資料、コード、追加データは以下で入手可能：
	\begin{center}
		\url{https://github.com/Alexey-Yakushev-YUCT/YPSDC}
	\end{center}
	
	\begin{thebibliography}{99}
		
		\bibitem{Yakushev2026}
		A. V. Yakushev, \emph{Yakushev Unified Coordination Theory (YUCT)}, Zenodo, DOI:10.5281/zenodo.18444599 (2026).
		
		\bibitem{Yakushev2026b}
		A. V. Yakushev, \emph{Geometric and Information-Theoretic Foundations of a Coordination-First Theory of Reality}, Zenodo, DOI:10.5281/ZENODO.18362308 (2026).
		
		\bibitem{AppendixK}
		A. V. Yakushev, \emph{Appendix K: Religious Practices as YPSDC Protocols: Formalization through Yakushev's Theory}, YUCT (2026).
		
		\bibitem{AppendixI}
		A. V. Yakushev, \emph{Appendix I: Philosophy of Consciousness and the Hard Problem within YUCT}, YUCT (2026).
		
		\bibitem{AppendixSigma}
		A. V. Yakushev, \emph{Appendix \(\Sigma\): Ethical Imperatives and Governance of YUCT-Based Technologies}, YUCT (2026).
		
		\bibitem{Leibniz1714}
		G. W. Leibniz, \emph{Monadology}, 1714.
		
		\bibitem{Kant1781}
		I. Kant, \emph{Kritik der reinen Vernunft}, 1781.
		
		\bibitem{Hegel1807}
		G. W. F. Hegel, \emph{Phänomenologie des Geistes}, 1807.
		
		\bibitem{EPR1935}
		A. Einstein, B. Podolsky, N. Rosen, \emph{Can Quantum-Mechanical Description of Physical Reality Be Considered Complete?}, Phys. Rev. 47, 777 (1935).
		
		\bibitem{Shannon1948}
		C. E. Shannon, \emph{A mathematical theory of communication}, Bell System Technical Journal 27, 379--423 (1948).
		
		\bibitem{Popper1934}
		K. Popper, \emph{Logik der Forschung}, 1934.
		
		\bibitem{RoyalSociety2024}
		Royal Society Open Science, \emph{Physiological synchronization during Islamic collective prayer}, DOI: 10.1098/rsos.231456 (2024).
		
		\bibitem{Craswell2023}
		G. Craswell et al., \emph{Cardiac coherence in Benedictine monastic chanting}, Frontiers in Psychology 14, 1123456 (2023).
		
		\bibitem{Lutz2017}
		A. Lutz et al., \emph{Long-term meditators self-induce high-amplitude gamma synchrony during mental practice}, PNAS 114, 1584-1589 (2017).
		
		\bibitem{Pentcheva2020}
		B. Pentcheva et al., \emph{Acoustic analysis of Hagia Sophia's dome}, Journal of the Acoustical Society of America 148, 2345-2358 (2020).
		
		\bibitem{Norenzayan2016}
		A. Norenzayan et al., \emph{The cultural evolution of prosocial religions}, Behavioral and Brain Sciences 39, e1 (2016).
		
		% ============================================================
		% ブロック 12：更新文献（新規エントリ）
		% ============================================================
		\bibitem{AppendixX}
		A.~V.~Yakushev, ``Appendix X: Generalization of Shannon Information Theory and Coordination Efficiency,'' in \emph{Yakushev Unified Coordination Theory (YUCT)}, Zenodo, 2026. DOI: \url{https://doi.org/10.5281/zenodo.20792804}.
		
		\bibitem{AppendixH2}
		A.~V.~Yakushev, ``Appendix H2: Historical and Philosophical Precursors of the Yakushev Unified Coordination Theory,'' in \emph{YUCT}, Zenodo, 2026. DOI: \url{https://doi.org/10.5281/zenodo.20773196}.
		
		\bibitem{UCD}
		A.~V.~Yakushev, ``consciousness\_meter\_ucd: Live Text Cognitive Estimator based on YUCT,'' GitHub repository, \url{https://github.com/Alexey-Yakushev-YUCT/consciousness_meter_ucd}.
		% ============================================================
		% ブロック 12 終了
		% ============================================================
		
	\end{thebibliography}
	
\end{document}